\documentclass[11pt,oneside]{article}
\usepackage[T1]{fontenc}
\usepackage{lmodern}
\usepackage{amsmath,amssymb,amsthm,mathtools}
\usepackage[margin=1in]{geometry}
\usepackage{microtype,booktabs,longtable,array,enumitem}
\usepackage{xcolor,xurl}
\usepackage[colorlinks=true,linkcolor=blue!40!black,citecolor=blue!40!black,urlcolor=blue!45!black]{hyperref}
\usepackage{fancyhdr}
\hypersetup{pdftitle={Mathematics in iPRA: finite demands, root supports, and norm certificates},
pdfauthor={Research report prepared for Vladimir Reshetnikov},
pdfsubject={Primitive recursive finite observations, aggregate witnesses, and superelliptic zero counts}}
\setlength{\headheight}{14pt}
\pagestyle{fancy}
\fancyhf{}
\fancyhead[L]{\small Mathematics in intuitionistic primitive recursive arithmetic}
\fancyhead[R]{\small 15 November 2026}
\fancyfoot[C]{\thepage}
\setlength{\emergencystretch}{2em}
\setlength{\parskip}{0.25em}
\setlist{itemsep=0.25em,topsep=0.5em}
\newtheorem{theorem}{Theorem}[section]
\newtheorem{lemma}[theorem]{Lemma}
\newtheorem{proposition}[theorem]{Proposition}
\newtheorem{corollary}[theorem]{Corollary}
\theoremstyle{definition}
\newtheorem{definition}[theorem]{Definition}
\newtheorem{example}[theorem]{Example}
\theoremstyle{remark}
\newtheorem{remark}[theorem]{Remark}
\newcommand{\ipra}{\ensuremath{\mathrm{iPRA}}}
\newcommand{\N}{\mathbb N}
\newcommand{\Z}{\mathbb Z}
\newcommand{\Q}{\mathbb Q}
\newcommand{\R}{\mathbb R}
\newcommand{\C}{\mathbb C}
\newcommand{\abs}[1]{\lvert#1\rvert}
\newcommand{\norm}[1]{\lVert#1\rVert}
\newcommand{\dotminus}{\mathbin{\dot{-}}}

\title{\textbf{Mathematics in iPRA}\\[0.55em]
\large Finite demands, root supports, and norm certificates}
\author{Research report prepared for Vladimir Reshetnikov}
\date{15 November 2026}

\begin{document}
\maketitle
\thispagestyle{empty}
\begin{abstract}
We extend the finite-observation approach in the ProveIt iPRA proof
collection in three directions. First, a backward-demand theorem permits
iteration of represented objects using only quantifier-free induction:
one fixes a final precision, computes predecessor demands, and proves a
bounded invariant. A related construction compiles a predetermined
parameter-changing recurrence into two ordinary primitive recursions.
Second, nonempty finite rational sets provide aggregate witnesses.
We prove finite cover calculus, a prescribed-chain point constructor,
and quantitative optimization. Applied to monic polynomials, bounded
rational deflation and a product estimate produce a uniform PR name of
the entire root support, including at collisions, together with explicit
precision and adequacy bounds. Third, multiplication matrices give norm
certificates for general monic plane equations. For
$y^m=f(x)$ with $\gcd(m,\deg f)=1$, a unique top determinant term proves
that the norm of $\sum_{i<m}A_i(x)y^i$ has degree exactly
$\max_i(m\deg A_i+i\deg f)$. A finite Vandermonde argument converts
this into a sharp zero-count and proper-subspace fiber bound.

The report gives complete human arguments, identifies the inherited
rational-descent lemma, and supplies concrete Lean definitions and a
dependency-ordered formalization plan. The scheduling and finite-cover
cores are proved in native Lean; the main analytic and algebraic
applications remain named formalization targets, not completed embedded
\texttt{Prf PRAAx} derivations. Exact arithmetic programs and validation
records are included. Singular Puiseux coefficient selection, general
Diophantine height cutoffs, and arbitrary rational quotient-cover
producers remain outside the proved conclusions.
\end{abstract}

\vspace{0.5em}
\noindent\textbf{Source baseline:}
\texttt{052fbaebfc955294b4c02f6532ab7e7a3f760f7a}.\par
\noindent\textbf{Scope:} a targeted extension of the current collection,
including its three existing framework reports; not a certification of
all 54 delivered packets.\par
\noindent\textbf{Formal status:} five new Lean modules compile with
Lean 4.32.0. Their proved lemmas and unproved proposition targets are
distinguished throughout.

\clearpage
\section*{Result guide}
\addcontentsline{toc}{section}{Result guide}

The general method is to ask what finite observation a later argument
needs, then construct and prove exactly that observation with a primitive
recursive resource bound. Infinite names and abstract mathematical
objects remain useful notation, but the induction takes place on finite
codes. Three cases require different treatments: demands that grow under
iteration, outputs whose identity is unstable under perturbation, and
finite certificates whose nonvanishing carries genuine mathematical
content.

\begin{center}
\begin{tabular}{@{}p{0.23\textwidth}p{0.66\textwidth}@{}}
\toprule
Main component & Concrete result\\
\midrule
Backward demands & From $C(n,D(n,k))\to C(n+1,k)$, compute
$h_N=K$, $h_t=D(t,h_{t+1})$ and derive $C(N,K)$ from
$C(0,h_0)$ by a guarded bounded invariant.\\[6pt]
Aggregate witnesses & Nonempty finite lists with
$H(S_s,S_t)\le2^{-s}+2^{-t}$ support PR distances, prescribed-chain
points, finite images, and approximate minimization.\\[6pt]
Root supports & For monic degree $n$, coefficient precision
$E(t)=n(t+2)+b$ suffices to produce lists with
$H(S_t,S_u)\le2^{-t-2}$ for $u\ge t$, once
$2^b\ge4nR^{n-1}$.\\[6pt]
Norm certificates & A computed nonzero norm bounds distinct common
zeros by its degree. For coprime superelliptic exponents the exact degree
is automatic, and every proper weighted-evaluation subspace has at most
$L$ points in a supplied distinct list.\\
\bottomrule
\end{tabular}
\end{center}

The root theorem uses the earlier FTA packet's explicit rational descent
lemma, whose article is included unchanged. Its new construction performs
all deflations on finite rational polynomials; it never invokes an
unclocked evaluator for arbitrary PR programs. The norm theorem is an
independent finite-algebra argument. It supplies a wider geometric
fiber-count contract for the Siegel strategy, while preserving the
distinction between cardinality and a global height bound.

\clearpage
{\footnotesize\setlength{\parskip}{0pt}\tableofcontents}
\clearpage
\section{The target theory and the embedding discipline}
\label{sec:framework}

\subsection{What is being embedded}

The target is the first-order intuitionistic primitive recursive arithmetic
of \path{ConciseIPRA/Lean/ConciseIPRA/PR.lean} \cite{repo-core}.
Its terms have a symbol for each \emph{specified} primitive recursive
function. These symbols are generated by zero, successor, projections,
composition, and primitive recursion. The nonlogical axioms are the
defining equations of these symbols, successor injectivity,
$Sx\ne0$, and induction for quantifier-free formulas. The logic is
intuitionistic natural deduction with equality, including the usual rules
for $\bot,\wedge,\vee,\to,\forall,\exists$. The quantifier-free induction
schema permits free numerical parameters:
\[
 A(0,\mathbf p)\ \wedge\
 \forall j\bigl(A(j,\mathbf p)\to A(j+1,\mathbf p)\bigr)
 \quad\longrightarrow\quad\forall j\,A(j,\mathbf p).
\]
Here $A$ must be quantifier-free. The theory does not acquire full
induction merely because every function appearing in a formula is PR.

We embed \emph{presentations, observations, and constructions}. A rational
number is an integer numerator and positive denominator; a Gaussian
rational is a pair of rationals; matrices and polynomials are finite
lists. Their equality tests and arithmetic are fixed PR algorithms.
Canonical reduction by gcd is convenient but optional: cross
multiplication gives an exact equality test on valid rational codes.
Malformed inputs receive fixed default outputs so that every numerical
function remains total.

One possible list coding is
\[
 \operatorname{nil}=0,\qquad
 \operatorname{cons}(a,b)=2^a(2b+1).
\]
For a positive code, a bounded search for its exponent of two recovers
the head, and integer division recovers the tail. The tail is smaller
than the code. Iterating for at most the original code therefore
implements length, lookup, and validation by PR functions. Nested lists,
signed integer pairs, and finite tables introduce no new sort into the
first-order theory. This inefficient coding suffices for the logical
claim; the reference implementations use ordinary Python and Lean lists.

An infinite object is described by a fixed PR family of finite
observations. For example, a real name $x(\mathbf p,k)\in\Q$ obeys
\[
 |x(\mathbf p,k)-x(\mathbf p,\ell)|\le2^{-k}+2^{-\ell}.
 \tag{N}\label{eq:regular-name}
\]
The displayed condition is an ordinary first-order formula after coding.
It is not a new real-number sort. Its intended real value is useful
semantic notation; every proof below ultimately concerns finitely many
rational inequalities. Equality of represented objects is generally
unbounded and is never assumed decidable.

There are two different kinds of uniformity. A theorem can be uniform in
the degree, dimension, finite coefficient list, and numerical parameters
of one fixed PR evaluator. Alternatively, it can be a metatheorem
instantiated separately for each PR evaluator. Neither permits one
unclocked PR interpreter for all PR program codes: if such an interpreter
$U$ existed, $n\mapsto U(n,n)+1$ would itself be PR, and its index
would contradict its own evaluation. A bounded universal machine with an
explicit clock is harmless. Our root construction uses fixed coefficient
evaluators and a single rational-polynomial algorithm; the norm
construction has only finite inputs.

\subsection{Why bounded arithmetic is enough for the finite parts}

\begin{lemma}[Bounded formulas have Boolean representatives]
\label{lem:booleanization}
For each fixed bounded numerical formula $B(\mathbf x)$ whose atoms
compare PR numerical terms, one can construct a PR Boolean function
$b_B$ such that
\[
 \ipra\ \vdash\ B(\mathbf x)\ \longleftrightarrow\ b_B(\mathbf x)=1.
\]
This is a metatheorem for fixed syntax, not an evaluator for arbitrary
formulas and arbitrary PR symbol codes.
\end{lemma}

\begin{proof}
Natural-number equality and order have PR characteristic functions;
their specifications follow by quantifier-free induction on the usual
recursive comparison functions. Rational comparisons reduce to integer
comparisons after multiplication by positive denominators. Boolean
negation, conjunction, and disjunction handle the corresponding
connectives. Since the already constructed Boolean atoms are decidable,
implication is represented by $\neg b\vee c$: the verification is
intuitionistic case analysis on a Boolean value, not excluded middle for
arbitrary propositions.

For $\exists y<t\,C(\mathbf x,y)$, take the finite Boolean disjunction
of $b_C(\mathbf x,0),\ldots,b_C(\mathbf x,t-1)$. For a bounded universal
quantifier, take the finite conjunction. The empty disjunction is false
and the empty conjunction true. These equivalences can be verified
without assuming bounded induction in advance. Write $A(i)$ for the
conjunction of the first $i$ Boolean values. From the bounded-universal
premise for a fixed terminal $T$, use the quantifier-free guarded
invariant $i\le T\to A(i)=1$. Conversely, fix an index $j$ and induct
on $i$ with the quantifier-free invariant
$j<i\wedge A(i)=1\to b_C(\mathbf x,j)=1$; split the successor case
into $j=i$ and $j<i$ by numerical decidability. After this induction,
introduce the bounded universal quantifier over $j$.

For the existential fold $E(i)$, maintain a first-success register
$w(i)$, with default $i$ when there has been no success. Its recurrence
either retains the old successful index or records $i$ on the next
successful test. The invariants
\[
 E(i)=1\ \to\ w(i)<i\wedge b_C(\mathbf x,w(i))=1,
 \qquad
 j<i\wedge b_C(\mathbf x,j)=1\ \to\ E(i)=1
\]
are quantifier-free, with $j$ a fixed parameter in the second induction.
They prove both implications and produce the existential witness.
Each fixed nesting of bounded quantifiers adds a fixed nesting of PR
recursions. Structural induction on the fixed formula completes the
construction. The detailed arithmetic groundwork is also provided by
the earlier finite-contract treatment \cite{repo-contracts}.
\end{proof}

In particular, induction on a bounded formula is an admissible derived
rule: replace it by its Boolean equation, use quantifier-free induction,
and translate back. To reason about a loop with terminal length $T$, use
the guarded invariant
\[
                         j\le T\ \longrightarrow\ I(j).
\]
Its Boolean representative is quantifier-free. For $j+1\le T$ the
local transition establishes the next invariant; when this guard fails,
the implication holds by decidable arithmetic. Thus no induction beyond
the terminal index is required.

The surrounding hypothesis context may contain unbounded validity
assertions, for example \eqref{eq:regular-name}. Such assertions remain
hypotheses. At a finite step one instantiates them at the finitely many
needed indices; one does not put them into the induction invariant.
If all those indices have a PR bound, their needed instances can instead
be collected into one bounded Boolean table. This is the precise role of
the observation-demand calculations that follow.

\subsection{Five separate obligations}

The previous frameworks already established bounded Booleanization,
guarded loops, certificate composition, finite observation contracts,
and several quantitative representation tools
\cite{repo-contracts,repo-embedding,repo-first-framework}. We retain that
foundation and make the following obligations explicit for each new
construction.
\begin{enumerate}
\item \textbf{Representation.} Specify the finite codes, validity
conditions, and what equality of the represented objects means.
\item \textbf{Production.} Give a fixed PR algorithm, including bounds
for every search and a default on invalid inputs. A PR checker alone
does not supply a passing certificate.
\item \textbf{Finite correctness.} Prove the output observations from
identified input observations. Supply PR demand bounds and bounded
induction invariants.
\item \textbf{Adequacy.} Explain exactly which mathematical conclusion
the observations express. Small residuals, a coherent witness, a finite
set of witnesses, and a complete list with a height bound are different
outputs.
\item \textbf{Formalization.} Construct the object-theory derivation,
or label the remaining derivation as a target. Native Lean correctness
and an embedded \texttt{Prf PRAAx} term are distinct achievements.
\end{enumerate}

There is no general rule that a true universal statement about PR
functions is provable in \ipra. The displayed proofs supply the needed
finite derivations; their eventual Lean compilation must reproduce
those derivations. Likewise, a complete analytic statement may still
require a new estimate. The framework organizes such an estimate and
its uses; it does not manufacture it from the name of a certificate.

\subsection{How the new components fit together}

Section~\ref{sec:backward-demands} handles a family of constructions whose
correctness at precision $k$ depends on an earlier stage at a different
precision. The proof fixes the final demand and computes all intermediate
demands backward, so its induction invariant remains bounded.

Section~\ref{sec:aggregates} supplies a finite-set interface when the task
requires coverage of possible witnesses but no canonical choice of one.
It gives arithmetic proofs of composition, distances, images, and
approximate optimization. Section~\ref{sec:rootsets} supplies a producer
for that interface from monic polynomial coefficient names, including
multiple roots.

Section~\ref{sec:norm-certificates} addresses a different obstacle:
turning a finite algebraic presentation into a certificate with enough
mathematical force to bound a fiber. A determinant can be constructed
uniformly, and its nonvanishing can be proved automatically for coprime
superelliptic exponents. This supplies a previously missing geometric
input to a Diophantine proof strategy; it does not supply that strategy's
remaining approximation theorem or global height bound.

\section{Backward demands replace induction on completed objects}
\label{sec:backward-demands}

A recurring obstruction in an arithmetic treatment of analysis is an
induction hypothesis of the form ``the object at stage $n$ is valid at every
precision.''  For example, validity of a real name is an unbounded universal
statement.  The same issue occurs when one proves all field laws at every
stage of a tower of algebraic extensions.  Quantifier-free induction does
not permit these statements as induction formulas.

The existing packets already provide bounded Booleanization, finite loops,
and composition of observation contracts
\cite{repo-contracts,repo-embedding}.  Moreover, the Platonic-solids packet
uses a backward sequence of coefficient budgets for its field-tower
argument \cite{repo-platonic}.  The contribution here is a single explicit
rule covering nonstationary and nonmonotone demands, a construction of
parameter-changing recursions, and a checked Lean realization of their
finite scheduling arguments.  This is a generalization and reusable
formalization of that method, not a claim that backward budgeting itself
is new.

\subsection{The exact finite transfer theorem}

All functions in this subsection are fixed primitive recursive functions;
the notation $\vec a$ denotes their ordinary numerical parameters.  Let
$C(n,k,\vec a)$ be a bounded formula: it describes a finite observation of
the object at stage $n$, with demand $k$.  It can, for example, assert all
the required inequalities with indices at most $k$.  It is not required
to be monotone in $k$.  Let $D(n,k,\vec a)$ be a demand transformer.  A local
step contract has the shape
\begin{equation}\label{eq:backward-local-step}
 C(n,D(n,k,\vec a),\vec a)
       \ \longrightarrow\ C(n+1,k,\vec a).
\end{equation}
The point of this contract is that its premise asks for one specified
finite observation, instead of asking that the predecessor object satisfy
its specification at every precision.

Fix a terminal stage $N$ and a terminal demand $K$.  Define the backward
array by ordinary primitive recursion:
\begin{align}
 b(0;N,K,\vec a)&=K,\nonumber\\
 b(j+1;N,K,\vec a)
   &=D\bigl(N\mathbin{\dotminus}(j+1),b(j;N,K,\vec a),\vec a\bigr),
   \label{eq:backward-array}\\
 h(t;N,K,\vec a)&=b(N\mathbin{\dotminus}t;N,K,\vec a).
   \label{eq:backward-horizon}
\end{align}
Here $\dotminus$ is truncated subtraction.  The values with $j>N$ or
$t>N$ merely make the displayed functions total.  The proof uses only
the intended range.  Direct substitution gives
\begin{equation}\label{eq:backward-schedule-laws}
 h(N)=K,\qquad h(t)=D(t,h(t+1),\vec a)\quad(t<N).
\end{equation}
For the second equation use
$N\dotminus t=(N\dotminus(t+1))+1$ and
$N\dotminus((N\dotminus(t+1))+1)=t$, valid when $t<N$.
These are elementary numerical equations proved by quantifier-free
induction; no backwards induction principle is being added.

\begin{theorem}[Backward-observation transfer]
\label{thm:backward-transfer}
Let $\Gamma(\vec a,N,K)$ be a finite context in which the stage variable
does not occur freely.  Suppose that \ipra\ derives
\[
 \Gamma\vdash C(0,h(0;N,K,\vec a),\vec a)
\]
and, with $n,r$ free,
\[
 \Gamma,\ n<N,\ C(n,D(n,r,\vec a),\vec a)
       \vdash C(n+1,r,\vec a).
\]
Then \ipra\ derives
\[
                 \Gamma\vdash C(N,K,\vec a).
\]
Only the single initial observation at demand $h(0)$ is required.  Neither
$D$ nor $C$ needs a monotonicity hypothesis.
\end{theorem}

\begin{proof}
Fix $\vec a,N,K$ throughout the argument.  Booleanize the fixed bounded
formula $C$: let $c(n,k,\vec a)$ be its primitive recursive characteristic
function, with an \ipra\ derivation of
$C(n,k,\vec a)\leftrightarrow c(n,k,\vec a)=1$.
Use the induction formula
\begin{equation}\label{eq:backward-invariant}
 J(t):\quad
 t\le N\ \longrightarrow\ c(t,h(t;N,K,\vec a),\vec a)=1.
\end{equation}
The numerical comparison $t\le N$ is also a primitive recursive test, so
$J$ is quantifier-free in the language of \ipra.

The base follows from the supplied initial derivation and $0\le N$.
For the successor step assume $J(t)$ and $t+1\le N$.  Then $t<N$ and
$t\le N$, hence the induction hypothesis gives $C(t,h(t),\vec a)$.
By \eqref{eq:backward-schedule-laws}, this is exactly
$C(t,D(t,h(t+1),\vec a),\vec a)$.  Instantiate the supplied local step at
$n=t$ and $r=h(t+1)$.  It gives $C(t+1,h(t+1),\vec a)$, and
Booleanization yields the consequent of $J(t+1)$.  This argument proves
the implication defining $J(t+1)$; it asserts nothing about out-of-range
stages.

Quantifier-free induction proves $\forall t\,J(t)$.  Instantiate at $N$,
use $h(N)=K$, and convert the Boolean equation back to $C(N,K,\vec a)$.
The context $\Gamma$ remains fixed during induction.  Its formulas need
not themselves be decidable or bounded.
\end{proof}

\begin{corollary}[Observational induction]
\label{cor:observational-induction}
Suppose the context $\Gamma(\vec a)$ is independent of $N,K$, \ipra\
derives $\Gamma\vdash C(0,k,\vec a)$ for arbitrary $k$, and it derives
\eqref{eq:backward-local-step} under $\Gamma$.  Then it derives
\[
       \Gamma\vdash\forall n\,\forall k\ C(n,k,\vec a).
\]
\end{corollary}

\begin{proof}
Fix $N,K$, substitute $h(0;N,K,\vec a)$ for $k$ in the base derivation,
and apply Theorem~\ref{thm:backward-transfer}.  Universal introduction over
$N,K$ gives the result.  The induction formula was
\eqref{eq:backward-invariant}; it was not $\forall k\,C(t,k,\vec a)$.
\end{proof}

For a constant precision loss $c$, namely $D(n,k)=k+c$, the schedule is
especially simple:
\begin{equation}\label{eq:constant-loss-budget}
                  h(t)=K+c(N\dotminus t).
\end{equation}
If $D(n,k)=k+d(n)$, then $h(0)=K+\sum_{n<N}d(n)$.  A nonlinear $D$ is
handled by the same primitive recursion \eqref{eq:backward-array}; no
closed formula or elementary growth bound is needed.  This remains a
fixed primitive recursive definition uniformly in the numerical inputs
$N,K,\vec a$.  It is not a universal evaluator for variable program syntax.

\subsection{Parameter-changing recursion without hidden strength}

Even defining an iterated represented object can conceal a difficulty.
The equations
\begin{equation}\label{eq:parameter-changing-recursion}
 x_0(k)=b_0(k),\qquad
 x_{n+1}(k)=F(n,k,x_n(D(n,k)))
\end{equation}
change the parameter $k$ in the recursive call.  Therefore one should not
simply declare that they are an instance of the primitive-recursion
constructor with $k$ held fixed.

\begin{theorem}[Compilation of a predetermined demand recurrence]
\label{thm:parameter-changing-recursion}
For fixed primitive recursive $b_0,D,F$, there is a fixed primitive
recursive function $X(N,K)$ satisfying
\eqref{eq:parameter-changing-recursion}.  Its equations have \ipra\ proofs.
\end{theorem}

\begin{proof}
Compute the schedule $h(t;N,K)$ by
\eqref{eq:backward-array}--\eqref{eq:backward-horizon}.  Then compute
\[
 v_0=b_0(h(0)),\qquad
 v_{t+1}=F(t,h(t+1),v_t),\qquad X(N,K)=v_N.
\]
This is a composition of two ordinary primitive recursions, with $N,K$
as fixed parameters in the second one.  All data are natural-number codes.

For $N=0$, $h(0;0,K)=K$, so $X(0,K)=b_0(K)$.  For the successor
equation compare the schedule for terminal data $(N+1,K)$ with that for
$(N,D(N,K))$.  Their demands at time $N$ are both $D(N,K)$.  The same
backward recurrence therefore gives identical demands at every time
$t\le N$.  More explicitly, quantifier-free induction on $j$ proves
\[
 b(j+1;N+1,K)=b(j;N,D(N,K));
\]
the equality of the relevant truncated subtractions supplies its step.
Substitution gives equality of the two schedules at $t\le N$.
Their initial forward values agree, and quantifier-free induction on
$t\le N$ gives equality of their forward values.  The last step of the
larger computation is consequently
\[
       X(N+1,K)=F(N,K,X(N,D(N,K))).
\]
Each induction here compares numerical terms; it does not quantify over
functions or assume an unbounded family of already constructed objects.
\end{proof}

The predetermined nature of $D$ matters.  The theorem does not license
arbitrary nested recursion in which a future demand is determined by a
not-yet-computed recursive output.  Such a definition needs a separate
primitive recursive construction or bound.  A bounded trace can support
adaptive queries, as in the earlier finite-contract frameworks, but
termination of an unbounded query process is not supplied by this theorem.

\subsection{Variable finite sets of demands}

A local mathematical argument often needs several predecessor facts.
Let $A(n,j,\vec a)$ be bounded, let $m(n,j,\vec a)$ be the number of
queries, and let $q(n,j,r,\vec a)$ be the $r$th query.  Suppose the local
derivation is
\begin{equation}\label{eq:multiple-demand-step}
 \left(\forall r<m(n,j,\vec a)\;
 A(n,q(n,j,r,\vec a),\vec a)\right)
       \longrightarrow A(n+1,j,\vec a).
\end{equation}
Define
\begin{align}
 C(n,H,\vec a)&:\Longleftrightarrow
       \forall j\le H\ A(n,j,\vec a),\label{eq:prefix-observations}\\
 D(n,H,\vec a)&=
       \max_{j\le H}\ \max_{r<m(n,j,\vec a)}q(n,j,r,\vec a).
       \label{eq:multiple-demand-envelope}
\end{align}
An empty maximum is zero.  Finite maxima are primitive recursive, even
when the number of queries varies with the input and none of their indices
is monotone.

\begin{corollary}[Finite-query observational induction]
\label{cor:finite-query-induction}
Given \ipra\ derivations of $A(0,j,\vec a)$ and
\eqref{eq:multiple-demand-step} under a fixed context, there is an \ipra\
derivation of $\forall n\,\forall j\,A(n,j,\vec a)$ under that context.
\end{corollary}

\begin{proof}
For $j\le H$ and $r<m(n,j,\vec a)$, the defining maximum gives
$q(n,j,r,\vec a)\le D(n,H,\vec a)$.  Thus $C(n,D(n,H,\vec a),\vec a)$
supplies every premise in \eqref{eq:multiple-demand-step}.  It yields
$A(n+1,j,\vec a)$ for each $j\le H$, hence $C(n+1,H,\vec a)$.
Apply Corollary~\ref{cor:observational-induction} to the prefix predicate.

There are two bounded conjunctions here, not a hidden universal invariant.
Concretely, if $a(n,j,\vec a)$ Booleanizes $A$, then $C(n,H,\vec a)$ is
Booleanized by the product of the $H+1$ Boolean values $a(n,j,\vec a)$.
The query premise is the analogous product over $r<m(n,j,\vec a)$.
Their prefix-fold equations and the finite-maximum inequalities have
quantifier-free induction proofs.  The proof remains within \ipra\ when
$H$ and the query counts are input-dependent.
\end{proof}

Multidimensional observations require no new logical rule.  Code a tuple
of indices by a fixed primitive recursive pairing function, and put
domain and range checks into its bounded predicate.  A query can carry a
tag naming which law, coordinate, or estimate is requested.  If a local
step uses several earlier stages, replace \eqref{eq:prefix-observations}
by the bounded history predicate
\[
             C(n,H):\Longleftrightarrow
                  \forall i\le n\;\forall j\le H\ A(i,j).
\]
Use a demand envelope at least $H$ and at least every requested query
index, and require each requested stage to be at most $n$.  Old history
entries persist because the envelope is at least $H$; the local step
establishes the new entries.  This proves the same transfer rule for a
finite directed dependency graph ordered by its stage index.  Cyclic
dependencies are not justified by this argument.

\subsection{Worked application: nonlinear iteration of a real name}
\label{subsec:nonlinear-name-iteration}

This example exhibits the computational content of the rule.  It avoids
both an induction on equality of real numbers and the assumption that
an arbitrary nested approximation program is primitive recursive.

Fix a rational $a$ with $0\le a\le4$ and a fixed primitive recursive
rational family $u(k)$.  Suppose its regularity observations are
\[
 \operatorname{Obs}_u(H):\Longleftrightarrow
  \forall i,j\le H\quad
       |u(i)-u(j)|\le2^{-i}+2^{-j},
\]
including positivity checks for all denominators used.  Put
\[
 p(r)=\min(1,\max(0,r)),\qquad T_a(r)=ar(1-r),
\]
and define the output family directly by the finite loop
\begin{equation}\label{eq:nonlinear-name-definition}
                v(n,k)=T_a^{\,n}\bigl(p(u(k+2n))\bigr).
\end{equation}
For each input $n,k$, this evaluates $u$ once and performs $n$ exact
rational updates.  Therefore it is a fixed primitive recursive family,
uniformly in $n,k,a$ and the numerical parameters of the fixed evaluator
$u$.  Rational arithmetic may be implemented by signed numerators and
positive denominators; reduction to lowest terms is unnecessary.

Three rational inequalities provide all local mathematics.  First,
$p$ takes values in $[0,1]$ and is nonexpansive:
$|p(r)-p(s)|\le|r-s|$.  To see the latter, assume $r\le s$ and split at
$0,1$.  On the middle interval $p$ is the identity; on either outer
interval it is constant.  Crossing an endpoint removes a nonnegative
part of the distance, so in every case
$0\le p(s)-p(r)\le s-r$.
Second, $T_a$ maps $[0,1]$ into itself, because
\[
       0\le r(1-r)\le\frac14
       \quad\text{follows from}\quad (2r-1)^2\ge0.
\]
Third, for $r,s\in[0,1]$,
\begin{equation}\label{eq:logistic-lipschitz}
 |T_a(r)-T_a(s)|
    =a|r-s|\,|1-r-s|\le4|r-s|.
\end{equation}
All cases and inequalities concern rational numbers and can be checked
after clearing positive denominators.

The finite iterate definition also gives the exact identity
\begin{equation}\label{eq:logistic-parameter-shift}
                     v(n+1,k)=T_a(v(n,k+2)).
\end{equation}
Indeed, both sides start from $p(u(k+2n+2))$ and apply the same map
$n+1$ times.  This identity can equivalently be obtained by
Theorem~\ref{thm:parameter-changing-recursion} with demand $D(n,k)=k+2$.

\begin{theorem}[Finite validity of the nonlinear iterates]
\label{thm:nonlinear-name-validity}
From $0\le a\le4$ and $\operatorname{Obs}_u(K+2N)$, \ipra\ proves
\[
 \begin{gathered}
      \forall k\le K\quad 0\le v(N,k)\le1,\\
      \forall k,\ell\le K\quad
       |v(N,k)-v(N,\ell)|\le2^{-k}+2^{-\ell}.
 \end{gathered}
\]
In particular, a globally regular input family gives a globally regular
output family at every numerical stage $N$.
\end{theorem}

\begin{proof}
Let $C(n,H)$ be the conjunction of the displayed range and regularity
assertions with $n,H$ in place of $N,K$.  It is bounded.  At $n=0$,
$v(0,k)=p(u(k))$; the range assertion and nonexpansiveness of $p$ give
$C(0,K+2N)$ from the input observation.

Assume $C(n,H+2)$.  For $k\le H$, the input $v(n,k+2)$ lies in $[0,1]$,
so \eqref{eq:logistic-parameter-shift} gives the next range assertion.
For $k,\ell\le H$, the same identity and
\eqref{eq:logistic-lipschitz} give
\[
 \begin{split}
 |v(n+1,k)-v(n+1,\ell)|
 &\le4|v(n,k+2)-v(n,\ell+2)|\\
 &\le4\bigl(2^{-(k+2)}+2^{-(\ell+2)}\bigr)\\
 &=2^{-k}+2^{-\ell}.
 \end{split}
\]
Thus $C(n,H+2)\to C(n+1,H)$ is the local finite contract.
The backward schedule is $h(t)=K+2(N-t)$ for $t\le N$.
Apply Theorem~\ref{thm:backward-transfer}.  Its induction formula is the
Boolean equation for $C(t,K+2(N-t))$, guarded by $t\le N$.
Only after this finite argument is complete do we universally quantify
over $N,K$ when the input is globally regular.
\end{proof}

In the ordinary external interpretation of regular names, if $u$ names
$x$, then $v(N,\cdot)$ names $T_a^N(p(x))$.  The same Lipschitz estimate
gives an error at most
$4^N2^{-(k+2N)}=2^{-k}$ at precision $k$.  When $x\in[0,1]$, this is
the usual nonlinear orbit $T_a^N(x)$.  No completed real-number sort or
completeness axiom occurs in the internal theorem.  Nor does the theorem
promise a fixed input precision adequate for arbitrarily many iterations:
the explicit cost is two additional bits per iteration.

\subsection{Lean status and the exact object-theory endpoint}
\label{subsec:backward-lean-status}

The file \path{lean/BackwardDemand.lean} imports only \texttt{Std} and
contains proved declarations for the following items:
\begin{itemize}
\item \texttt{horizon\_step} and \texttt{horizon\_coherence}: the schedule
  equations and their compatibility when one final step is added;
\item \texttt{compiledRecursion\_succ}: the successor equation of the
  two-loop implementation of parameter-changing recursion;
\item \texttt{transfer} and \texttt{observational\_induction}: finite
  transfer and subsequent universal closure;
\item \texttt{finite\_query\_induction}: reduction of a variable finite
  family of predecessor queries to a prefix maximum;
\item \texttt{constant\_shift\_transfer}: the explicit accumulated budget
  $K+Nc$.
\end{itemize}
All these declarations compiled with Lean 4.32.0.  Their axiom audit lists
only \texttt{propext} and \texttt{Quot.sound}; it lists neither
\texttt{sorryAx} nor \texttt{Classical.choice}.  This checks the native
Lean scheduling arguments.  The Lean predicates in these declarations
are arbitrary \texttt{Prop}-valued functions; Lean does not infer from
their types that they are primitive recursive or bounded.

The corresponding object-theory target is the repository judgment
$\texttt{Prf PRAAx}\;\Gamma\;\varphi$ \cite{repo-core}.
A complete formalization has the following concrete steps.
\begin{enumerate}
\item Supply fixed \texttt{PR} syntax for $D$, the test Booleanizing $C$,
  truncated subtraction, and the two recursions.  The recursion argument
  of \texttt{PR.prec} comes first; use projections and composition to put
  $N,K,\vec a$ in the parameter slots.  The formulas
  \eqref{eq:backward-array}--\eqref{eq:backward-horizon} are the complete
  numerical definitions needed for this step.
\item Prove the subtraction identities and
  \eqref{eq:backward-schedule-laws} as object equalities, using
  \texttt{PRA.precZero}, \texttt{PRA.precSucc}, and equality congruence.
  The schedule-coherence proof is another quantifier-free equality
  induction.  No evaluator over arbitrary \texttt{PR} syntax is required.
\item Supply the actual object derivations of bounded Booleanization and
  the application's base and local-step contracts.  For finite queries,
  use the explicit Boolean products in
  Corollary~\ref{cor:finite-query-induction} and the primitive recursive
  prefix maximum in \eqref{eq:multiple-demand-envelope}.
\item Form the syntax of \eqref{eq:backward-invariant}.  Its base has type
  $\Gamma\vdash\varphi[0/t]$.  Its step has type
  $\varphi::\operatorname{shift}(\Gamma)\vdash\varphi[S(t)/t]$.
  Use substitution, \texttt{Prf.eqSubst1}, and the supplied local contract
  to construct this step.  These shifted contexts are necessary to avoid
  capturing the arbitrary endpoint parameters.
\item Apply \texttt{PRA.qfInd} to that quantifier-free formula, then apply
  universal elimination at the term $N$.  Rewrite $h(N)=K$ and use the
  reverse Booleanization implication.  Universal introduction yields
  the corollary only when its variables are absent from the fixed context.
\item For the nonlinear application, introduce signed-rational code
  operations, clipping, and exact iteration; prove the three displayed
  rational inequalities and the parameter-shift identity.  Instantiate
  the preceding derivation transformer with their bounded tests.
\end{enumerate}
The native Lean file completes the scheduling proofs, including the
nonmonotone and variable-query cases.  The PR syntax compiler, general
Booleanization compiler, and the final object derivations of the nonlinear
example remain explicit formalization tasks.  A semantic theorem proved
with native Lean induction is not being advertised as an already
constructed \ipra\ proof object.

Finally, the rule is intentionally sufficient rather than an unrestricted
proof-extraction principle.  A local derivation and a fixed PR demand must
be supplied.  Finite consistency of a table does not prove global validity
of a proposed infinite object.  A loop that repeatedly asks whether a
represented real is nonzero has no guaranteed terminating demand merely
because every finite query is legitimate.  The diagonal obstruction to a
universal unclocked PR evaluator also remains intact
\cite{repo-contracts,repo-embedding}.

\section{Finite aggregates instead of premature witness selection}
\label{sec:aggregates}

\subsection{A representation with finite, decidable observations}

Work first in $\Q^d$ with the sup metric
$d(x,y)=\max_{i<d}|x_i-y_i|$. This metric is rational-valued and PR on
finite vector codes. Everything below also applies to any fixed coded
rational metric with proved triangle and symmetry laws. The complex
root application uses squared Euclidean comparisons; its estimates
imply the sup-metric estimates because $\|v\|_\infty\le|v|_2$.

For nonempty finite lists $S,T$, repetitions allowed, put
\[
 h(S,T)=\max_{x\in S}\min_{y\in T}d(x,y),\qquad
 H(S,T)=\max\bigl(h(S,T),h(T,S)\bigr).
\]
Minima and maxima mean finite scans with the first index breaking a tie.
All comparisons are rational. Equivalently, $H(S,T)\le r$ says
\[
 (\forall x\in S\ \exists y\in T\ d(x,y)\le r)
 \ \wedge\
 (\forall y\in T\ \exists x\in S\ d(x,y)\le r).
 \tag{H}\label{eq:finite-cover}
\]
This is a bounded formula, hence has a Boolean checker by
Lemma~\ref{lem:booleanization}. Its existential quantifiers are finite
searches, not choices of infinite point names.

\begin{lemma}[Finite aggregate calculus]\label{lem:aggregate-calculus}
The following statements have finite proofs on the supplied lists.
\begin{enumerate}
\item $H(S,U)\le H(S,T)+H(T,U)$.
\item $H(S\mathbin{\|}S',T\mathbin{\|}T')\le
 \max(H(S,T),H(S',T'))$, where $\|$ concatenates lists.
\item With the sup product metric,
 $H(S\times S',T\times T')\le\max(H(S,T),H(S',T'))$.
\item If $d'(f(x),f(y))\le Ld(x,y)$ on the relevant finite inputs,
 then $H(f[S],f[T])\le LH(S,T)$.
\item Writing $d(z,S)=\min_{x\in S}d(z,x)$, one has
 $|d(z,S)-d(w,T)|\le d(z,w)+H(S,T)$.
\end{enumerate}
\end{lemma}

\begin{proof}
For the first assertion, choose for each supplied $x\in S$ a nearest
$y\in T$, then a nearest $z\in U$. The triangle inequality gives
$d(x,z)\le H(S,T)+H(T,U)$. The reverse direction is identical with the
lists reversed. These choices are least-index finite scans.

For concatenation, inspect which part of the concatenated list contains
the current index and use that part's cover. For products, choose the
nearest candidate in each factor and use the sup product metric. For
images, apply the assumed finite Lipschitz inequality to every selected
cover pair. Finally choose a minimizing $y\in T$ for $d(w,T)$ and an
$x\in S$ within $H(S,T)$ of it. Then
$d(z,S)\le d(z,w)+d(w,T)+H(S,T)$. Interchanging the pairs gives the
absolute-value bound. All statements about all list elements are bounded
conjunctions; prefix induction verifies their construction.
\end{proof}

\begin{definition}[Regular aggregate name]\label{def:aggregate-name}
A regular aggregate name is a fixed PR family $k\mapsto S_k$ of
nonempty finite rational lists satisfying
\[
                      H(S_s,S_t)\le e_s+e_t,
             \qquad e_j=2^{-j}.                 \tag{A}\label{eq:aggregate-name}
\]
The family can depend on numerical parameters. Its represented object
is the completion of these finite observations, not a decision procedure
for membership in a finite list of exact points.
\end{definition}

For a fixed rational point $z$, Lemma~\ref{lem:aggregate-calculus}(5)
shows that $k\mapsto d(z,S_k)$ is a regular rational name. Denote its
represented value by $D_S(z)$. If $z$ is itself a regular point name,
use $d(z_{k+1},S_{k+1})$: its change between indices $k,\ell$ is at
most $e_k+e_\ell$. This defines the same distance extension. In semantic
notation,
\[
 |D_S(z)-d(z,S_s)|\le e_s,\qquad
 |D_S(z)-D_S(w)|\le d(z,w).                       \tag{D}\label{eq:distance-name}
\]
The first inequality follows by comparing to $S_t$ and taking an
arbitrarily fine \emph{explicit} later observation. For a finite desired
slack $2^{-k}$, it is enough to take $t\ge k+1$ and a point precision
at least $k+1$. Limits here are shorthand for those finite inequalities.

One useful equality convention for aggregate names is
\[
                     S\simeq T\quad\Longleftrightarrow\quad
              \forall k\ H(S_k,T_k)\le2e_k.     \tag{E}\label{eq:aggregate-equality}
\]
It is an equivalence relation. Only transitivity needs explanation.
If $S\simeq T$ and $T\simeq U$, then for every $v\ge k$,
\[
 H(S_k,U_k)\le2e_k+6e_v.
\]
The left side and $2e_k$ are rationals. If their difference were positive,
its numerator and denominator give an explicit binary-length bound for
$v$ with $6e_v$ smaller than that difference, a contradiction. Thus
$H(S_k,U_k)\le2e_k$. This uses rational decidability, not Markov's
principle. Equivalent names give identical distance names by the same
finite comparison. Equality \eqref{eq:aggregate-equality} is itself
unbounded and is not a Boolean test on program indices.

\subsection{Realizing a point by a prescribed chain}

A fixed PR sequence of finite lists gives more than an arbitrary
qualitative assertion of approximability. It permits the following
explicit point constructor.

\begin{theorem}[Least-neighbor chain]\label{thm:aggregate-chain}
Given a regular aggregate name, an index $s$, and an entry $q\in S_s$,
one can construct a PR point name $x^{s,q}$ such that
\[
                  D_S(x^{s,q})=0,
                  \qquad d(x^{s,q},q)\le3\,2^{-s}.
\]
The construction is uniform in $s$, the entry index, and the numerical
parameters of the fixed list evaluator. It does not assert extensional
or continuous selection of a preferred point from the aggregate.
\end{theorem}

\begin{proof}
Put $q_s=q$. Once $q_j\in S_j$ has been produced, let $q_{j+1}$ be
the first nearest entry of $S_{j+1}$. Nonemptiness and
\eqref{eq:aggregate-name} give
\[
 d(q_j,q_{j+1})\le e_j+e_{j+1}.
\]
For $s\le j<v$, a finite telescoping sum gives
\[
 d(q_j,q_v)\le\sum_{r=j}^{v-1}(e_r+e_{r+1})
            =3e_j(1-2^{-(v-j)})\le3e_j.          \tag{C}\label{eq:chain-tail}
\]
This is induction on a finite table of rational points. Its invariant is
membership of the last point in the last finite list, together with the
displayed sum bound; it does not assert the existence of a limiting point.

Define the output at precision $t$ to be $x^{s,q}(t)=q_{s+t+2}$.
Equation~\eqref{eq:chain-tail} implies, for $u\ge t$,
\[
 d(x^{s,q}(t),x^{s,q}(u))\le3\,2^{-s-t-2}\le2^{-t}.
\]
Thus the output is a regular PR point name. The same finite bound gives
its distance from $q_s$ at most $3e_s$. Membership has a particularly
direct finite verification. Its distance-name sample at index $k$ is
\[
 d\bigl(x^{s,q}(k+1),S_{k+1}\bigr)
 =d(q_{s+k+3},S_{k+1})
 \le e_{s+k+3}+e_{k+1}\le\tfrac58 e_k,
\]
because $q_{s+k+3}\in S_{s+k+3}$. These nonnegative rational samples
therefore name zero. The largest list index needed in this check is
$s+k+3$; all evaluation indices and loop lengths are explicit.

Every step calls the same fixed PR list evaluator and performs a finite
minimum scan. Iteration is primitive recursion on $j$. If an invalid
input produces an empty list, choose the zero vector as a default; the
proof above concerns the stated valid-input hypotheses.
\end{proof}

Define membership in the represented support by $D_S(x)=0$. The theorem
shows that each finite approximant lies within $3e_s$ of a represented
member. Conversely, \eqref{eq:distance-name} shows that every member has
distance at most $e_s$ from $S_s$. This is a statement about its minimum
distance name, not a selection of an exact-radius witness for an arbitrary
real point. Replacing $S_s$ by $S_{s+2}$ absorbs the factor
three if a unit-error convention is desired. This semantic explanation
uses a prescribed PR sequence; it is not an invocation of countable
choice for arbitrary abstract approximation data. Richman's completed
finite-set approach is a useful broader antecedent \cite{richman2000}.

\subsection{Images and optimization, with their exact output meaning}

Finite union and product of regular names remain regular by
Lemma~\ref{lem:aggregate-calculus}. For an approximate represented map,
suppose a fixed PR evaluator $f(q,p)$ obeys the numerical contract
\[
 d'\bigl(f(q,p),f(r,p')\bigr)
          \le Ld(q,r)+2^{-p}+2^{-p'}              \tag{L}\label{eq:approx-lipschitz}
\]
on all rational inputs, with a given nonnegative rational $L$.
A specified invariant domain containing every sample and intermediate
point suffices instead if its membership is included in the contracts.
Choose a natural number $c$ with $2^c\ge L+1$, using integer bit lengths
after clearing the denominator, and put
\[
                  T_k=\bigl[f(q,k+c):q\in S_{k+c}\bigr].
\]
The finite covers and \eqref{eq:approx-lipschitz} give
\[
 H(T_s,T_t)\le(L+1)(e_{s+c}+e_{t+c})\le e_s+e_t.
\]
This proves a PR constructor for the completed finite-image object.
The finite observation indices are exactly the displayed ones. If $S\simeq S'$, the same calculation gives $H(T_k,T'_k)\le2(L+1)e_{k+c}\le2e_k$, so the finite-image constructor is extensional on aggregate names.

The action on represented point names is also explicit:
\[
             \widehat f(x)(k)=f(x(k+c),k+c).
\]
Contract \eqref{eq:approx-lipschitz} and regularity of $x$ bound its
sample difference by $(L+1)(e_{s+c}+e_{t+c})\le e_s+e_t$.
Equivalent point names give equivalent outputs by the same estimate,
and $d'(\widehat f(x),\widehat f(y))\le Ld(x,y)$ in represented
notation. Thus later occurrences of $f(x)$ refer to this prescribed
extension, rather than to a separately assumed completed function.

The word ``image'' does not conceal a preimage selection principle.
If $D_S(x)=0$, then $D_T(f(x))=0$ by the Lipschitz contract.
We have not proved a uniform constructor of $x$ from an arbitrary named
$y$ with $D_T(y)=0$. Likewise, the distance for a union is
$\min(D_S,D_T)$; a proof that it vanishes is not a decision of the
disjunction $D_S=0\vee D_T=0$. This distinction lets aggregate
representations preserve useful information without reintroducing the
selection difficulty they were meant to avoid.

\begin{theorem}[Value and approximate witnesses for optimization]
\label{thm:aggregate-optimization}
Let $f$ be a real-valued map satisfying \eqref{eq:approx-lipschitz}.
Its infimum over the represented support has a PR regular name.
For each $k$, one can produce a PR point name $x^{(k)}$ in that support
whose value is within $2^{-k}$ above the infimum. No one exact minimizing
point is selected by this theorem. The same assertions hold for the
supremum by replacing $f$ with $-f$.
\end{theorem}

\begin{proof}
Let $a_s=\min_{q\in S_s}f(q,s)$. Pick finite minimizers and use the
two directed covers. Contract \eqref{eq:approx-lipschitz} gives
\[
                         |a_s-a_t|\le(L+1)(e_s+e_t).
\]
Thus $a_{k+c}$, with the same $c$ as above, is a regular PR name;
write its value as $a$. In particular $|a-a_s|\le(L+1)e_s$.
If $x$ belongs to the support, choose a nearest finite point to a
sufficiently accurate rational observation of $x$. Equations
\eqref{eq:distance-name} and \eqref{eq:approx-lipschitz} imply
$a\le f(x)$ after making their explicit errors arbitrarily small.
For a desired slack $\varepsilon$, it suffices to choose $s$ with
$(4L+2)e_s<\varepsilon$ and point precision at least $s$.
This proves the lower-bound part directly from finite observations.

Choose $q_s\in S_s$ with $f(q_s,s)=a_s$ by a finite scan. The chain
theorem gives $x_s$ in the support with $d(x_s,q_s)\le3e_s$.
Writing $f(q_s)$ for its represented value, the same contract yields
\[
 f(x_s)\le f(q_s)+3Le_s
          \le a_s+(3L+1)e_s
          \le a+(4L+2)e_s.
\]
Choose $b$ with $2^b\ge4L+2$ and use $s=k+b$. Then $x^{(k)}=x_s$
has the required error. Both the finite minimizing index and the
subsequent point-name evaluator are PR. The approximating minimizers can
switch components as $k$ changes; nothing in this proof asserts their
coherence as a single minimizing name.
\end{proof}

\subsection{A reusable application to spectra}

The root-support producer in Section~\ref{sec:rootsets} combines with
this calculus to give an aggregate spectrum for an $n$ by $n$ matrix
of PR complex names, $n\ge1$. This statement concerns the characteristic
polynomial and its root support; it does not require eigenvectors,
rank decisions, or a Jordan normal form.

Here is an explicit coefficient precision budget. Use complex entry
names regular in the Euclidean norm. Suppose every sampled
entry has modulus at most an integer $B\ge1$; one can take
$B=3+\lceil\max_{i,j}\|a_{ij}(0)\|_1\rceil$.
For names given with a sup-coordinate convention, sample one extra
precision bit to obtain the Euclidean convention. A characteristic
coefficient is a signed sum of principal minors. There are at most
$2^n$ such minors in any coefficient and at most $n!$ determinant terms
per minor. A product of $r\le n$ entries changes by at most
$rB^{r-1}\delta$ when corresponding entries change by at most $\delta$;
expand the product difference one factor at a time. Thus the safe bound
\[
 A=2^n n!(B+1)^n,\qquad
 C=2^n n!\,n(B+1)^{n-1}
\]
respectively bounds every lower coefficient and its sensitivity to the
maximum entry error. Compute the characteristic coefficients by finite
determinant sums, using entry precision $s+c$ with $2^c\ge C$.
Their Cauchy error is then at most $e_s+e_t$. All loops and factorials
are PR uniformly in $n$.

Apply the root theorem, then the maximum construction to $z\mapsto|z|$.
Rational square-root bisection supplies approximations to $|z|$, with
the inequality $||z|-|w||\le|z-w|_2\le2\|z-w\|_\infty$;
thus use the safe Lipschitz constant $L=2$ in the sup-metric calculus.
The relevant disk bound is explicit.
This gives a PR name for the spectral radius, and finite coverage of
every spectral point, including at repeated eigenvalues. The result is
a direct application of our proved interfaces, not a claim to have
formalized the spectral theorem. The coefficient calculations are
compatible with the generic-matrix approach of the Cayley--Hamilton
packet \cite{repo-cayley}.

\subsection{Lean realization of the finite interface}

\path{lean/FiniteHausdorff.lean} defines directed covers, a finite
Boolean checker, the two-sided relation, and regular aggregate names.
It proves checker reflection, relational cover composition, finite union,
finite image transport, and the rational-metric triangle theorem.
Its exact axiom audit is included: the purely relational composition
lemma uses no axioms, while the rational triangle instance inherits
\texttt{Classical.choice} from Lean's standard rational-order lemmas.
This is a native-library dependency, not an added object-iPRA axiom.
The human finite rational proof above and its future arithmetization
remain independent obligations.

The companion \path{lean/AggregateTargets.lean} supplies concrete
rational-vector distance, finite minimum, least-neighbor iteration,
distance-name membership, represented-map evaluation, and approximate
minimizer programs. Their full chain, value, and witness theorems are
named unproved propositions. These definitions specify the remaining
mathematical statements rather than assuming them as structure fields.

To finish the object proof, encode the metric and lists, establish the
arithmetic triangle law by rational cross multiplication, prove the
finite-minimum specification by a prefix invariant, and replace every
list quantifier in \eqref{eq:finite-cover} by its Boolean fold. For the
chain theorem, the invariant is exactly membership plus
\eqref{eq:chain-tail}. For optimization, use the finite minimum and
the displayed $(L+1)$ and $(4L+2)$ budgets. No compactness theorem from
an ambient classical real analysis library is needed.

\section{An aggregate replacement for singular root selection}
\label{sec:rootsets}

The Puiseux packet identifies a real obstruction, rather than a missing
estimate. For \(Y^2-X^2-a\), an exact Puiseux coefficient at the returned
ramification index distinguishes \(a=0\) from \(a>0\). Such a constructor
cannot be uniform in arbitrary fast coefficient names. The two existing
framework packets therefore construct simple branches with a supplied
derivative gap. Here we address a different output: the finite set of
roots, observed at a specified precision \cite{repo-puiseux,repo-contracts}.
Multiple roots and collisions
are allowed. We do not recover a singular Puiseux expansion.

The use of completed finite sets and multisets in constructive algebra
has an established antecedent in Richman's \emph{The fundamental theorem
of algebra: a constructive development without choice} (2000)
\cite{richman2000}. The
contribution here is the explicit first-order primitive recursive
construction, its finite precision budget, and its application to the
representation boundary recorded in this repository. In particular, the
set representation below deliberately forgets multiplicity; a theorem
about matching multisets would be a different result.

\subsection{The finite metric relation}

Write \(\mathbb Q[i]\) for the Gaussian rationals. For \(z=x+iy\), put
\[
 N(z)=x^2+y^2,\qquad \|z\|_1=|x|+|y|.
\]
The notation \(|z|\le r\), for rational \(r\ge0\), means the rational
comparison \(N(z)\le r^2\). Euclidean norms used in intermediate estimates
abbreviate their elementary squared-norm consequences; they are not
additional input objects. We use
\(N(zw)=N(z)N(w)\), \(|z|\le\|z\|_1\le2|z|\), and finite
triangle and Cauchy--Schwarz inequalities.

For nonempty finite lists \(S,T\) and rational \(r\ge0\), define
\[
 \operatorname{Cover}(S,T,r)
 \;:\Longleftrightarrow\;
 \forall z\in S\ \exists w\in T\; N(z-w)\le r^2.
\]
Put \(\operatorname{Close}(S,T,r)\) for the conjunction of both directed
covers. This is a bounded decidable predicate. Repeated list entries are
permitted and have no effect on the predicate. Thus this is precisely a
finite presentation of Hausdorff closeness of the underlying supports.

\begin{lemma}[Finite product stability]\label{lem:product-stability}
Let \(S=(z_1,\ldots,z_n)\) and \(T=(w_1,\ldots,w_n)\), where \(n\ge1\),
and suppose every entry has norm at most the rational \(R\ge1\). Write
\[
 F(X)=\prod_{j=1}^n(X-z_j),\qquad
 G(X)=\prod_{j=1}^n(X-w_j).
\]
Suppose corresponding coefficients below degree \(n\) differ in norm
by at most \(\delta\), where \(\delta\ge0\). Set
\(C=nR^{n-1}\). If \(C\delta\le r^n\), then
\(\operatorname{Close}(S,T,r)\).
\end{lemma}

\begin{proof}
For \(z\in S\), finite polynomial evaluation gives
\[
 |G(z)|=|G(z)-F(z)|
 \le\delta\sum_{i=0}^{n-1}|z|^i
 \le C\delta\le r^n.
\]
But \(G(z)=\prod_j(z-w_j)\). Test the finitely many rational comparisons
\(N(z-w_j)\le r^2\). If all fail, each squared distance is strictly
greater than \(r^2\), and multiplication yields
\(N(G(z))>r^{2n}\), a contradiction. Bounded search therefore produces
the required index. Exchange the lists for the reverse cover. The
argument also works when \(r=0\): a finite product of nonnegative
rationals can vanish only if a factor vanishes. Every search and
induction in this proof is on explicit finite data.
\end{proof}

No assignment pairing all roots injectively was used. For example,
\((0,0,1)\) and \((0,1,1)\) have identical supports, although their
multiplicities differ. The lemma must not be read as a multiplicity
matching theorem.

\subsection{A finite approximate-root dependency}

We reuse one explicit human-proof result from
\path{algebra/FTA_iPRA/FTA_iPRA.tex}, Section~4, ``Uniform descent
for a rational polynomial'' \cite{repo-fta}. For a nonconstant Gaussian-rational
polynomial of degree at most \(d\ge1\), its primitive recursive
\(\operatorname{Descent}\) operation satisfies
\[
 |P(z+\Delta)|\le q_d|P(z)|,
 \qquad q_d=1-\frac1{4D_d},\qquad D_d=2^d4^{2d^2}.
 \tag{RD}\label{eq:rational-descent}
\]
The supplied argument uses a bounded dyadic radius search, a dominant
term at three consecutive scales, and a bounded rational direction
calculation. It does not assume the existence of a root. This is the
only nontrivial algebraic-existence dependency of the construction below.
It has a written proof in the repository; we do not claim its complete
\texttt{Prf PRAAx} formalization is available.

Here is the precise finite corollary we require, with its proof. Given
\(0<\varepsilon\le1\), let
\[
 B=\max(1,\|P(0)\|_1),\quad
 K=\lceil B/\varepsilon\rceil,\quad J=4D_dK.
\]
Start at zero and perform exactly \(J\) descent steps, assigning the
zero step when the residual is already zero. Finite Bernoulli arithmetic
gives \(q_d^{4D_d}\le1/2\). Hence the output \(z\) satisfies
\[
 |P(z)|\le Bq_d^J\le B2^{-K}\le\varepsilon,
\]
because \(2^K\ge K+1\ge B/\varepsilon\). This defines a single
primitive recursive operation on rational polynomial codes, uniformly in
the degree and requested precision. Only rational polynomials occur;
there is no evaluation of an arbitrary primitive recursive program code.

For a slower but simpler executable interface, the same corollary yields
a completely bounded grid search. If \(P\) is monic and its lower
coefficients have norm at most an integer \(A\ge1\), any point with
residual at most one has norm at most \(R=A+2\). Indeed, for
\(u=|z|>A+2\), the finite geometric identity gives
\[
 |P(z)|\ge u^d-A\frac{u^d-1}{u-1}>1.
\]
Set \(U=2R+1\) and
\(L_P=\max(1,\sum_{i=1}^d i\|p_i\|_1U^{i-1})\).
Take \(m=\lceil4L_P/\varepsilon\rceil\). Use the corollary at accuracy
\(\varepsilon/2\), and round its output to the square grid of spacing
\(1/m\) in \([-R,R]^2\). The distance moved is at most \(1/m\), and
the finite power-difference identity bounds the residual change by
\(L_P/m\le\varepsilon/4\). Thus the grid contains a point of residual
at most \(\varepsilon\). Scanning its \((2Rm+1)^2\) points is another
uniform PR approximate-root algorithm. This proof of grid completeness
depends on (RD), not on a classical compactness or FTA principle.

\subsection{Uniform finite factorization without nested coefficient names}

For a polynomial with Gaussian-rational coefficients, let
\(\|P\|_{\mathrm{coef},1}=\sum_i\|p_i\|_1\).
This norm is rational, subadditive, and submultiplicative, as follows
directly by expanding the finite convolution.

\begin{theorem}[Uniform rational approximate factorization]
\label{thm:approx-factor}
Let \(P\in\mathbb Q[i][X]\) be monic of degree \(n\ge1\), and let
\(0<\eta\le1\) be rational. From these finite inputs a primitive
recursive operation returns \(z_1,\ldots,z_n\in\mathbb Q[i]\) such that
\[
 \left\|P-\prod_{j=1}^n(X-z_j)\right\|_{\mathrm{coef},1}\le\eta.
\]
The operation is uniform in \(n\). Its mathematical proof uses (RD).
\end{theorem}

\begin{proof}
Choose an integer \(A_0\ge1\) bounding the \(\ell^1\) norm of every
coefficient of \(P\). Before performing any deflation, compute
\[
 A_{j+1}=(A_j+1)(2A_j+5)^n\quad(0\le j<n),\qquad
 L=\prod_{j=0}^{n-1}(2A_j+5),\qquad
 \varepsilon=\frac{\eta}{2nL}.
 \tag{AF}\label{eq:factor-budget}
\]
These are finite integer recursions. Let \(P_0=P\). At stage \(j<n\),
\(P_j\) is monic of degree \(n-j\). Apply the approximate-root
operation at tolerance \(\varepsilon\) and obtain \(z_{j+1}\). Perform
synthetic division, an exact finite recurrence, to obtain
\[
 P_j(X)=(X-z_{j+1})P_{j+1}(X)+r_j,
 \qquad r_j=P_j(z_{j+1}).
 \tag{SD}\label{eq:synthetic-division}
\]
The quotient is monic of degree \(n-j-1\).

Inductively, every coefficient of \(P_j\) has \(\ell^1\) norm at most
\(A_j\). The just-proved residual radius bound gives
\(|z_{j+1}|\le A_j+2\), because \(\varepsilon\le1\). Therefore
\(\|z_{j+1}\|_1\le2A_j+4\). The synthetic recurrence, beginning with
leading quotient coefficient one, bounds each quotient coefficient by
\[
 (A_j+1)(1+\|z_{j+1}\|_1)^n
 \le (A_j+1)(2A_j+5)^n=A_{j+1}.
\]
For clarity, the elementary recurrence estimate is obtained by repeatedly
expanding \(b_{k-1}=p_k+z_{j+1}b_k\): its right side is a sum of at
most \(n\) terms, each dominated by a term of
\((A_j+1)(1+\|z_{j+1}\|_1)^n\). Thus it does not require a loop
invariant about infinitely many coefficients.

After all \(n\) divisions, \(P_n=1\). Substituting
\eqref{eq:synthetic-division} repeatedly gives the exact identity
\[
 P(X)-\prod_{k=1}^n(X-z_k)
   =\sum_{j=0}^{n-1}r_j\prod_{k=1}^{j}(X-z_k).
 \tag{TE}\label{eq:factor-telescoping}
\]
Here the empty product is one. Since \(\|r_j\|_1\le2\varepsilon\),
submultiplicativity gives
\[
 \left\|P-\prod_{k=1}^n(X-z_k)\right\|_{\mathrm{coef},1}
 \le 2\varepsilon\sum_{j=0}^{n-1}
       \prod_{k=0}^{j-1}(2A_k+5)
 \le2\varepsilon nL=\eta.
\]
Each division manipulates a finite rational coefficient list; the outer
loop has exactly \(n\) stages. Consequently this construction avoids the
higher-type iteration issue in repeatedly deflating exact Cauchy names.
Its cost is large but primitive recursive.
\end{proof}

\subsection{A uniform finite-set name for all roots}

Fix one primitive recursive Gaussian-rational coefficient evaluator
\(a(i,s,\mathbf p)\). Parameters \(\mathbf p\), degree \(n\ge1\), and
precision indices are ordinary numerical inputs. Assume, for \(i<n\),
\[
 |a(i,s,\mathbf p)-a(i,t,\mathbf p)|\le2^{-s}+2^{-t}.
 \tag{CN}\label{eq:coefficient-name-root}
\]
The leading coefficient is exactly one. Suppress \(\mathbf p\) below.
Let
\[
 A=3+\left\lceil\max_{i<n}\|a(i,0)\|_1\right\rceil,
 \quad R=A+2,\quad C=nR^{n-1}.
\]
Then every sampled coefficient has Euclidean norm at most \(A-1\).
In the approximate-factorization procedure one can use \(A_0=2A\),
which bounds every sampled coefficient's \(\ell^1\) norm.
Choose an integer \(b\) with \(2^b\ge4C\), by bounded binary-length
calculation, and set
\[
 E(t)=n(t+2)+b,\qquad
 P_t(X)=X^n+\sum_{i<n}a(i,E(t))X^i.
\]
Apply Theorem~\ref{thm:approx-factor} to \(P_t\) with
\(\eta_t=2^{-E(t)}\). Denote its output list by \(S_t\) and its exact
product polynomial by \(F_t\).

\begin{theorem}[Collision-tolerant root-support names]
\label{thm:root-support}
The lists \(S_t\) form a primitive recursive finite-set name satisfying
\[
 u\ge t\quad\Longrightarrow\quad
 \operatorname{Close}(S_t,S_u,2^{-t-2}).
 \tag{RS}\label{eq:root-support-fast}
\]
All entries of all lists have norm at most \(R\). The construction is
uniform in the numerical degree and parameters of the fixed evaluator.
Its coefficient approximation and finite Hausdorff assertions have
\(\ipra\) proofs relative to the written descent lemma (RD).
\end{theorem}

\begin{proof}
Each lower coefficient of \(F_t\) has norm at most \(A\), because its
distance from the sampled coefficient is at most \(\eta_t\le1\).
Every entry of \(S_t\) is an exact rational root of the monic rational
polynomial \(F_t\). The radius argument above, now with residual zero,
therefore gives the common bound \(R=A+2\). This bound comes from the
\emph{final product coefficients}; we do not confuse it with the much
larger temporary bounds \(A_j+2\) used during deflation.

For \(u\ge t\), corresponding coefficients of \(F_t,F_u\) differ by
at most
\[
 2^{-E(t)}+(2^{-E(t)}+2^{-E(u)})+2^{-E(u)}
 \le4\cdot2^{-E(t)}.
\]
Consequently
\[
 C\,4\cdot2^{-E(t)}
 \le2^{-n(t+2)}=(2^{-t-2})^n.
\]
Lemma~\ref{lem:product-stability} proves \eqref{eq:root-support-fast}.
No equality or separation decision about represented coefficients or
roots occurs anywhere in the construction.
\end{proof}

The exact retained information is a nonempty finite-set name, together
with monic product polynomials converging coefficientwise to the input.
The next proposition proves that its zero-distance locus is exactly the
root locus, with an explicit finite-precision argument.

\begin{proposition}[Adequacy of the root-support representation]
\label{prop:root-adequacy}
Let \(z\) be a represented point, and supply an integer bound
\(Z\ge\max(R,1)\) with \(|z|\le Z\). Put
\[
 D_t(z)=\min_{w\in S_t}|z-w|,\quad
 C_Z=nZ^{n-1},\quad
 M_Z=A\sum_{i=1}^n iZ^{i-1}.
\]
The distance is a represented nonnegative real, computable by finite
minimum and square-root operations with prescribed precision. Define
\[
 \operatorname{Member}(z,S)\;:\Longleftrightarrow\;
 \forall k\ \exists t\ge k\;D_t(z)\le2^{-k}.
\]
Then \(\operatorname{Member}(z,S)\) holds if and only if \(P(z)=0\).
The forward production of indices from a root is primitive recursive.
Moreover the following error bounds hold:
\begin{align}
 |F_t(z)-P(z)|&\le2C_Z2^{-E(t)},\label{eq:root-eval-error}\\
 |P(z)|&\le M_ZD_t(z)+2C_R2^{-E(t)}.
 \label{eq:root-distance-residual}
\end{align}
\end{proposition}

\begin{proof}
Each lower coefficient of \(F_t\) differs from the represented input
coefficient by at most \(2\cdot2^{-E(t)}\): one contribution is the
approximate-factorization error and one is the sampled-name error.
The finite sum of powers at \(|z|\le Z\) proves
\eqref{eq:root-eval-error}. The power-difference identity proves
\(|P(z)-P(w)|\le M_Z|z-w|\) on this disk. At \(w\in S_t\),
\(F_t(w)=0\), and \(|w|\le R\), so
\(|P(w)|\le2C_R2^{-E(t)}\). This proves the inequality in
\eqref{eq:root-distance-residual} for each individual distance
\(|z-w|\), hence also for their finite minimum. No exact minimizing
index for represented distances has to be selected. For example a
finite minimum of real names can be implemented with
\(\min(u,v)=(u+v-|u-v|)/2\), with each arithmetic precision prescribed.

For completeness, here is a fully finite precision budget in the
root-to-membership direction. Let \(z(s)\) be its rational samples,
with sample error at most \(2^{-s}\) in represented notation. Set
\[
 U=Z+1,\quad C_U=nU^{n-1},\quad
 M_U=A\sum_{i=1}^n iU^{i-1},\quad T=C_U+M_U,
\]
and let \(\ell(v)\) be any primitive recursive integer choice with
\(2^{\ell(v)}\ge\max(1,v)\), for example binary length after rounding
\(v\) upwards. The rational polynomial evaluations
\[
 V_s=z(s)^n+\sum_{i<n}a(i,s)z(s)^i
\]
have error at most \(T2^{-s}\) relative to \(P(z)\); this follows by
separately varying the coefficients and the point. Thus \(P(z)=0\)
gives \(|V_s|\le T2^{-s}\). This is an explicit normalized
vanishing-error formulation of the root equation, not an invocation
of a limit principle. Direct coefficient comparison now gives the
entirely finite inequality
\[
 |F_t(z(s))|
 \le 2C_U2^{-E(t)}+(C_U+T)2^{-s}.
 \tag{FE}\label{eq:finite-root-evaluation}
\]
Given \(k\), choose
\[
 t=k+1+\ell(2C_U),\qquad
 s=\max\{k+1,\ n(k+3)+\ell(C_U+T)\}.
\]
The first summand in \eqref{eq:finite-root-evaluation} is at most
\(2^{-n(k+3)}\), because
\(E(t)\ge n(k+3)+\ell(2C_U)\); the second has the same bound.
Their sum is at most \(2^{-n(k+1)}\). All points in this last
calculation are rational. Testing their finitely many squared
distances and applying the product argument produces \(w\in S_t\)
with \(|z(s)-w|\le2^{-k-1}\). Since \(s\ge k+1\), adding the
sample error gives \(|z-w|\le2^{-k}\), as required. Both indices
and the successful finite search are fixed PR constructions. In a
pure sample formulation, the displayed represented inequalities are
obtained by substituting these precisions and the ordinary rational
slack-removal lemma; no real trichotomy is used.

Conversely, assume membership. Given a residual target \(2^{-k}\),
take \(j=k+\ell(M_Z+2C_R)\), and instantiate membership to obtain
\(t\ge j\) with \(D_t(z)\le2^{-j}\). Since \(E(t)\ge t\),
\eqref{eq:root-distance-residual} yields
\[
 |P(z)|\le(M_Z+2C_R)2^{-j}\le2^{-k}.
\]
This proves the root equation at every requested precision. The
existential index is used by intuitionistic existential elimination;
no unbounded search for it is assumed.
\end{proof}

By \eqref{eq:root-support-fast}, the distances \(D_t(z)\) have the
same explicit Cauchy error bound as the sets. Accordingly the displayed
membership formula is equivalent to zero distance from the represented
set. The generic finite-set chain construction supplies actual PR point
names when such a realization is needed. No arbitrary infinite choice
of labels is implicit in the root-support theorem.

\subsection{What the new representation resolves}

Lemma~\ref{lem:product-stability} is also a finite perturbation theorem.
For two monic input polynomials whose sampled coefficient differences
are at most \(\delta\), with splitting errors \(\eta,\eta'\), lists
in a common radius \(R\) are \(r\)-close whenever
\[
 nR^{n-1}(\delta+\eta+\eta')\le r^n.
 \tag{CP}\label{eq:coefficient-perturbation}
\]
This gives an explicit H\"older precision budget even at repeated roots.
For coefficients depending on another represented variable, combine a
supplied coefficient continuity modulus with
\eqref{eq:coefficient-perturbation}; this produces a modulus for the
root-support map. Neither a simple-root derivative bound nor root
separation is required.

For example, the obstruction family \(Y^2-X^2-a\) has a uniformly
represented root support at every represented input pair \((X,a)\).
For real \(X\), \(a\ge0\), the familiar semantic description is
\(\{-\sqrt{X^2+a},\sqrt{X^2+a}\}\); at \((0,0)\) the support is the
singleton \(\{0\}\). Coefficient perturbations of order \(\varepsilon^2\)
suffice for root-support accuracy of order \(\varepsilon\). The theorem
never has to decide whether the two roots coincide. It therefore repairs
the output contract for tasks needing distances, images, approximate
optimization, or coverage of all roots. It does not provide individually
accessible singular Taylor/Puiseux coefficients, multiplicity matching,
or a continuously labeled root on a parameter domain.

\subsection{Formalization obligations}

The file \path{lean/RootSetTargets.lean} defines Gaussian-rational
arithmetic, finite products, coefficient error bounds, the bounded cover
predicate, and the complete finite approximate-root constructor from a
supplied descent map with the displayed \(q_d,J\) budget. It also defines
point-name evaluation, normalized residual equality, and membership by
rational sample coverage. The forward adequacy indices are actual
definitions. Six named propositions specify the remaining correctness
proofs, including the descent-to-root bridge and the adequacy equivalence.
These propositions are \emph{unproved targets}, not axioms. The intended order
is: (1) polynomial convolution and synthetic-division identities;
(2) finite squared-norm inequalities and product stability;
(3) the explicit (RD) proof from the FTA packet;
(4) bounded approximate-root production; (5) the quantifier-free
deflation invariant and telescoping error bound; and (6) the fixed-name
precision substitution proving \eqref{eq:root-support-fast}.

For an object-theory proof, encode lists and Gaussian rationals as natural
numbers. Replace every list quantifier by its finite Boolean fold, retain
\eqref{eq:coefficient-name-root} as a hypothesis instantiated only at
the finitely needed indices, and induct on the finite deflation table.
The invariant includes coefficient bounds, degree, monicity, the exact
partial telescoping identity, and the accumulated rational error bound.
An ordinary Lean theorem over \texttt{Rat} is a useful preliminary layer;
the final goal is a derivation in the embedded \(\ipra\) calculus.
The reference program checks finite certificates with exact fractions.
It does not purport to execute the enormous worst-case descent schedule
or to verify the complete object-theory derivation.

\section{Automatic norm certificates and sharp finite zero counts}
\label{sec:norm-certificates}

The earlier embedding packet proves a sharp zero bound for
$y^2=f(x)$ when $\deg f$ is odd, using
$A^2-fB^2$. It also gives a general finite-algebra spanning-certificate
interface. The remaining question is how to \emph{construct} useful
certificates for wider classes without an unspecified elimination
algorithm \cite{repo-embedding,repo-siegel}. The present section supplies such a constructor for every
monic equation in one coordinate, and an automatic nonvanishing proof
for the coprime superelliptic family. For this family it replaces the
divisor-based zero-count obligation in the Siegel packet by explicit
polynomial and matrix arithmetic.

All fields used below have exact, decidable equality and primitive
recursive arithmetic on finite codes, with their field laws available
in the object theory. Rational numbers are the principal example. A
fixed explicitly presented number field works as well once those
arithmetic contracts have been proved. Arbitrary Cauchy names for real
or complex coefficients do not provide this interface. Lists of points
are always finite input data; no set of all points is formed.

\subsection{The determinant constructor for a monic plane equation}

Let $K$ be such a coded field and let
\[
 F(X,Y)=Y^m+\sum_{i<m}f_i(X)Y^i,\qquad m\geq1.
\]
Every polynomial in $Y$ has a remainder of degree less than $m$ after
division by $F$, computed by monic division. Write
$G=\sum_{i<m}A_i(X)Y^i$ for the remainder of the function whose zeros
are to be counted. For each $j<m$, reduce $GY^j$ modulo $F$ and put
its coefficient vector in column $j$ of an $m$ by $m$ matrix $M_G(X)$.
Define
\[
                         N_G(X)=\det M_G(X).
\]
The matrix and determinant are finite polynomial computations. In
particular this definition does not presume an algebraic closure,
factorization, a divisor, or a terminating search for a certificate.

\begin{theorem}[Finite norm zero bound]\label{thm:general-norm}
If $N_G\ne0$, every finite list of pairwise distinct points
$(x,y)\in K^2$ satisfying $F(x,y)=G(x,y)=0$ has length at most
$\deg N_G$. Nonvanishing of the computed polynomial is a finite
coefficient test. The statement includes singular curves and
ramified fibers.
\end{theorem}

\begin{proof}
First recall the elementary finite root bound: a nonzero polynomial of
degree $r$ over a field has at most $r$ distinct roots in a supplied
list. Synthetic division removes $Y-b$ at a listed root $b$. At a
different listed root $c$, its nonzero value $c-b$ implies that the
quotient still vanishes. A finite induction removes the product of
the distinct linear factors, whose degree cannot exceed $r$. This
is a proof by polynomial operations on the input list.

Group the supplied common zeros by their $x$-coordinates. Fix one
group with common abscissa $a$ and distinct ordinates
$b_1,\ldots,b_t$. The preceding root bound, applied to the monic
polynomial $F(a,Y)$, gives $t\leq m$. Set
\[
               e(b)=(1,b,\ldots,b^{m-1}).
\]
The definition of $M_G$ and evaluation of the monic-division
identities give
\[
 e(b_\ell)M_G(a)=G(a,b_\ell)e(b_\ell)=0
                    \qquad(1\leq\ell\leq t).
\]
Extend these $t$ rows to a square matrix $U$ by adjoining the
standard rows with indices $t,t+1,\ldots,m-1$, using zero-based
column indices. Its upper-left $t$ by $t$ block is the Vandermonde
matrix on the distinct $b_\ell$, and its lower-left block is zero
with lower-right block the identity. Thus
\[
             \det U=\prod_{1\leq i<j\leq t}(b_j-b_i)\ne0.
\]
The Vandermonde identity itself has a finite inductive proof.
Replace the last ordinate by an indeterminate $T$. The determinant
has degree at most $t-1$ in $T$ and vanishes at the preceding
$t-1$ distinct ordinates because two rows then agree. The root
removal argument already proved above divides out all
$T-b_i$. Its leading coefficient is the preceding $(t-1)$ by
$(t-1)$ Vandermonde determinant, by expansion in the last row.
Induction, starting with the determinant one in size one, proves
the formula after substitution $T=b_t$.

Every entry in the first $t$ rows of $UM_G(X)$ vanishes at $X=a$.
Synthetic division therefore removes a factor $X-a$ from each such
entry. Multilinearity in the rows shows
\[
                        (X-a)^t\mid\det(UM_G(X)).
\]
Since $\det U$ is a nonzero scalar and
$\det(UM_G)=\det U\,N_G$, it follows that
$(X-a)^t\mid N_G$. This step counts the \emph{number of distinct
points} in the fiber, not merely the number of abscissae. In
particular a ramified point with $b=0$ causes no division by $b$;
all divisions in $U^{-1}$, if one chooses to compute it, are by the
witnessed nonzero differences $b_j-b_i$ and by a nonzero determinant.
The proof does not need to compute that inverse.

For different abscissae, multiply these divisibilities. Here is an
explicit cancellation argument that avoids a factorization oracle.
Suppose $N_G=CQ$ after removing the factors belonging to previous
abscissae. At a new abscissa $a$ we have $C(a)\ne0$. From
$(X-a)^t\mid CQ$ and $t>0$, evaluation gives $Q(a)=0$. Write
$Q=(X-a)Q_1$ by synthetic division and cancel one $X-a$ from the
two polynomial expressions. Repeat exactly $t$ times. This removes
the new factor without disturbing the previous ones.

Consequently, if the groups have sizes $t_a$, then
\[
               \prod_a(X-a)^{t_a}\mid N_G.
\]
The product is monic of degree $\sum_a t_a$, which equals the
original list length. The assumption $N_G\ne0$ now proves the
required bound.
\end{proof}

There is also a fully explicit coarse degree estimate. Suppose all
nonzero $A_i$ have degree at most $a$ and all nonzero $f_i$ have
degree at most $c$, where $a,c\geq0$. The multiplication matrix for
$Y$ has entries of degree at most $c$: its nontrivial last column
is $(-f_0,\ldots,-f_{m-1})^T$. Therefore
$M_G=\sum_{i<m}A_iT^i$ has entries of degree at most
$a+(m-1)c$. Expanding its determinant gives
\[
              \deg N_G\leq m\bigl(a+(m-1)c\bigr)
                       \quad\text{if }N_G\ne0.
\]
The nonvanishing check is essential for a general monic equation.
For example $F=Y^2-X^2$ and $G=Y-X$ give $N_G=0$ and infinitely
many common rational zeros. The next theorem supplies a structural
nonvanishing proof for an important uniform class.

\subsection{Coprime superelliptic equations: no nonvanishing oracle}

Fix $m\geq2$ and a nonzero $f\in K[X]$ with positive degree $d$.
Assume
\[
                             \gcd(m,d)=1.
\]
For $G=\sum_{i<m}A_i(X)Y^i\ne0$, define its weighted degree by
\[
 W(G)=\max_{\substack{i<m\\A_i\ne0}}
                   \bigl(m\deg A_i+d i\bigr).
\]
The maximum has a unique maximizing index $k$. Indeed equality of
two displayed weights implies $m\mid d(i-j)$, hence $m\mid i-j$,
and $0\leq i,j<m$ then forces $i=j$. Coprime cancellation here can
be derived from the integer Bezout identity produced by bounded
Euclidean arithmetic.

\begin{theorem}[Exact degree of the superelliptic norm]
\label{thm:norm-degree}
For $F=Y^m-f(X)$, the determinant constructor above gives
\[
                       N_G\ne0,\qquad \deg N_G=W(G).
\]
If $k$ is the unique maximizing index, its leading coefficient is
\[
 \operatorname{sgn}(\sigma_k)
       \operatorname{lc}(A_k)^m\operatorname{lc}(f)^k,
 \qquad \sigma_k(j)=(j+k)\bmod m.
\]
\end{theorem}

\begin{proof}
Use row and column indices $r,j\in\{0,\ldots,m-1\}$. There is
exactly one $i$ in this range with $r\equiv i+j\pmod m$.
Writing $i+j=r+mq$, we have $q\in\{0,1\}$, and reduction of
$Y^{i+j}$ gives the particularly simple matrix entry
\[
                       (M_G)_{rj}=A_i f^q.
\]
Whenever this entry is nonzero,
\[
       m\deg (M_G)_{rj}+d r-d j
            =m\deg A_i+d i\leq W(G).
\]
Consider a nonzero determinant term associated to a permutation
$\sigma$. Sum this equality over its $m$ columns. The row and
column correction terms cancel because $\sigma$ permutes the
indices, so
\[
 m\deg\left(\prod_{j<m}(M_G)_{\sigma(j),j}\right)
       =\sum_{j<m}\bigl(m\deg A_{i_j}+d i_j\bigr)
       \leq mW(G).
\]
Equality is possible only if every $i_j=k$, by uniqueness of the
maximizing index. Exactly one permutation does this, namely
$\sigma_k(j)=(j+k)\bmod m$. Among these entries precisely $k$
wrap around the end of the index range. Their product is therefore
$A_k^m f^k$, of degree $m\deg A_k+dk=W(G)$.

Every other nonzero determinant term has strictly smaller degree.
Thus the top coefficient of the unique term cannot cancel. Its
value is the nonzero coefficient displayed in the statement, which
proves both assertions. This is a direct determinant computation;
no irreducibility theorem, valuation extension, or algebraic closure
is used.
\end{proof}

\begin{corollary}[Sharp weighted zero count]\label{cor:norm-zeros}
Every finite list of pairwise distinct points of $y^m=f(x)$ on
which $G(x,y)=0$ has length at most $W(G)$.
\end{corollary}
\begin{proof}
Apply Theorems~\ref{thm:general-norm} and~\ref{thm:norm-degree}.
\end{proof}

The degree hypothesis genuinely matters to the automatic proof.
When $m=d=2$, $f=X^2$ and $G=Y-X$, both summands have weight two;
their norm cancels identically. The general determinant theorem
remains valid whenever its computed norm is nonzero, even when
$\gcd(m,d)>1$.

The weighted bound is attained over $\Q$. Choose distinct rational
numbers $a_1,\ldots,a_d$ and put
\[
           f(X)=1+\prod_{j=1}^d(X-a_j),\qquad G=Y-1.
\]
For $\gcd(m,d)=1$, the common zeros are exactly
$(a_1,1),\ldots,(a_d,1)$, and $W(G)=d$. For example $m=3,d=4$
and $(a_j)=(0,1,2,4)$ give
$f=X^4-7X^3+14X^2-8X+1$. This example illustrates sharpness of
the arithmetic zero bound; smoothness is neither required nor
claimed by the theorem.

\subsection{A uniform proper-subspace fiber theorem}

For any $L\in\N$, let
\[
 \mathcal B_L=\{X^jY^i:0\leq i<m,\ j\geq0,\ mj+di\leq L\},
 \qquad n_L=|\mathcal B_L|.
\]
The constant monomial belongs to this list, so $n_L\geq1$.
Enumerate its entries lexicographically, and define
$\Phi_L(x,y)\in K^{n_L}$ by evaluating them in that order.
The list is obtained by the loops $i<m$ and $j\leq L$ followed by
a numerical test; no symbolic basis search occurs.

\begin{corollary}[Subspace fibers for superelliptic curves]
\label{cor:norm-subspace}
For every proper subspace $V\subsetneq K^{n_L}$ supplied by a
finite matrix of spanning rows of rank less than $n_L$, every
finite list of distinct points on $y^m=f(x)$ whose evaluation
vectors lie in $V$ has length at most $L$.
\end{corollary}
\begin{proof}
Bounded Gaussian elimination produces a nonzero annihilating
linear functional: choose a free column, assign its unknown one,
set the other free unknowns to zero, and solve the pivot equations
backwards. Its evaluation on $\Phi_L$ is a nonzero canonical
polynomial $G=\sum_{i<m}A_i(X)Y^i$ with $W(G)\leq L$.
Every point under consideration is a zero of $G$.
Corollary~\ref{cor:norm-zeros} gives the result.
\end{proof}

If a separately proved finite-tuple theorem covers the evaluation
vectors of each supplied list of points in a decidable class $S$
by at most $M$ such proper subspaces, then every distinct list
from $S$ has length at most $ML$. Assign each point to the first
covering subspace by bounded exact membership search; each of the
at most $M$ groups has size at most $L$. This last finite counting
step is the existing Siegel interface. The new contribution is an
automatic, sharp proof of its geometric fiber contract for all
coprime exponents $m,d$, as well as a computable norm test for a
general monic plane equation.

This result does not construct the quantitative approximation
cover, a complete list of rational or integral points, or a global
height cutoff. It therefore does not prove general Siegel or
Faltings. Nor does it impose smoothness or genus hypotheses that
would be needed for those larger theorems. It resolves the stated
finite zero-count dependency on a substantially wider class than
the previous odd-degree hyperelliptic case.

\subsection{Arithmetic compilation and Lean obligations}

For rational coefficients, use canonical signed numerators and
positive denominators, coefficient lists for univariate polynomials,
and lists of such lists for $G$. Monic reduction runs for the
input $Y$-degree. A companion matrix, finite matrix powers, and a
determinant expansion construct the norm. Permutations can be
enumerated by scanning all length-$m$ lists of indices less than
$m$ and retaining those with no repeated entries: the candidate
count is $m^m$. Thus even the deliberately inefficient determinant
algorithm is one fixed primitive recursive program, uniformly in
$m$. Euclidean and Gaussian elimination loops are bounded by
their input lengths or matrix dimensions. Exact coefficient size
may grow, but integer and rational arithmetic on codes remains
primitive recursive.

The human proof translates by quantifier-free loop invariants.
Monic division tracks the coefficient identity
$H=QF+R$ and the degree bound on its current remainder. Matrix
construction tracks, for every bounded column index, the
coefficient identity for $GY^j$ modulo $F$. The determinant
degree argument uses finite products, bounded permutation lists,
and a Boolean test for the unique cyclic shift. Grouping points
tracks that the accumulated groups partition the input list;
root removal tracks $N_G=CQ$ and the accumulated degree. All
coefficient equalities, bounded membership assertions, and
nonzero tests are Boolean folds on these finite data. The field
laws and the arbitrary supplied list remain external parameters
of these inductions. No induction formula asserts a universal
claim about all points of the curve.

The file \path{lean/NormTargets.lean} gives executable
coefficient-list definitions, a general monic-plane companion matrix,
monic reduction, the superelliptic matrix, determinants, weighted
degrees, finite-point checks, and named propositions for the main
theorems. It also defines weighted monomial evaluation vectors, row-span
witnesses, and nonzero annihilator certificates for the subspace result.
These propositions are
\emph{unproved targets}; compiling their definitions proves none
of their mathematical content. A complete formalization should
proceed in the following dependency order:
\begin{enumerate}
\item Prove normalization, convolution, evaluation, degree of
  nonzero products, and synthetic division for the chosen codes.
\item Prove the monic-reduction identities and the evaluation-row
  equation for the multiplication matrix.
\item Prove determinant multiplicativity, row multilinearity,
  the Vandermonde identity, and finite root removal; assemble
  Theorem~\ref{thm:general-norm}.
\item Prove the unique-weight lemma by Bezout cancellation,
  the entry weight equality, and the unique cyclic-shift term;
  assemble Theorem~\ref{thm:norm-degree}.
\item Prove the finite Gaussian annihilator construction and
  use the existing finite-cover counting argument.
\item Reify the fixed arithmetic programs and each displayed
  bounded invariant in the repository's primitive-recursive
  syntax, and build the corresponding embedded
  \texttt{Prf PRAAx} derivations. Native Lean matrix theorems alone
  are not this final object-theory step.
\end{enumerate}

The reference program \path{code/superelliptic_norm.py}
constructs these matrices and determinants with exact rational
polynomial arithmetic. It checks leading degrees against the
formula, determinant evaluation identities, the hyperelliptic
specialization, a sharp cubic-cover example, and repeated factors
coming from two-point and ramified fibers. These are finite
implementation checks supporting the definitions, separate from
the all-input human proofs above.

\section{A concrete path to complete Lean formalization}
\label{sec:formalization}

\subsection{Three levels that must remain distinct}

The report contains complete human arguments for the stated finite and
represented-output theorems. Their iPRA accounting identifies fixed PR
programs, bounded invariants, finite search limits, and the required
instances of unbounded input hypotheses. The root existence constructor
reuses the explicitly identified rational-descent lemma from the FTA
packet; an unchanged copy of its article is included in the archive.

The delivered Lean files are at two different implementation levels.
\path{BackwardDemand.lean} and \path{FiniteHausdorff.lean} prove
native Lean lemmas. \path{AggregateTargets.lean},
\path{RootSetTargets.lean}, and \path{NormTargets.lean} define
concrete algorithms and named propositions for their missing proofs.
A declaration of the form \texttt{def TheoremTarget : Prop := ...}
neither proves its proposition nor introduces it as an axiom.
All five files compile. The archive contains no admitted proof of a
main theorem and does not claim that these definitions constitute one.

The final logical-strength endpoint is different again:
\[
                  \texttt{Prf PRAAx [] sentence}.
\]
This type denotes a derivation in the repository's embedded arithmetic.
An arbitrary native Lean induction on \texttt{Prop} is not an instance
of the embedded theory's quantifier-free induction. In particular,
the native scheduling theorem deliberately accepts arbitrary predicates;
its iPRA application requires the separate boundedness and PR hypotheses
proved in Section~\ref{sec:backward-demands}.

\subsection{Definitions and theorem boundaries in the delivered modules}

\begin{longtable}{@{}p{0.24\textwidth}p{0.39\textwidth}p{0.30\textwidth}@{}}
\toprule
Module & Concrete content & Proof status\\
\midrule
\endhead
\texttt{BackwardDemand} & Backward schedule; forward evaluation;
parameter-changing recursion; finite query envelopes. & Schedule,
compilation equation, transfer, and finite-query theorems proved in
native Lean.\\[4pt]
\texttt{FiniteHausdorff} & Directed covers; Boolean checker; rational
metric; two-sided closeness; regular aggregate names. & Reflection,
composition, image, union, and triangle lemmas proved in native Lean.\\[4pt]
\texttt{AggregateTargets} & Rational vectors, exact minima, least-neighbor
chain, sampled membership, represented evaluation, and minimizers. &
Chain, distance, value, and approximate-witness goals are unproved
propositions.\\[4pt]
\texttt{RootSetTargets} & Gaussian rationals, polynomial operations,
descent clock, deflation, budgets, sampled root support and adequacy. &
Six targets: descent-to-root, product stability, factorization,
coherence, and two adequacy goals. All unproved.\\[4pt]
\texttt{NormTargets} & Rational polynomial lists, companion matrices, monic
reduction, determinant, norm, weighted degrees, points and subspaces. &
General and superelliptic norm, multiplicity, zero-count, and subspace
goals are unproved. Finite examples execute.\\
\bottomrule
\end{longtable}

\texttt{AggregateTargets.MemberZero} tests the rational samples of the
distance name at precision $k+1$ against $2^{-k}$. The least-neighbor
proof supplies the stronger $\tfrac58\,2^{-k}$ bound. Its optimization
targets specify both the value name and the complete two-parameter
point-name constructor. Thus an exact minimizer is not silently stored
as an assumed field of a structure.

\texttt{RootSetTargets} defines the supplied descent map's specification
in squared-norm arithmetic and implements its $J$-step iterator. The
remaining descent assumption has an exact upstream proof location.
Its adequacy targets use rational samples and bounded squared-distance
covers, avoiding comparisons of exact complex roots. Its integer
precision constants can be more generous than those in the paper:
choosing $b=4C$ instead of a binary length still guarantees $2^b\ge4C$.
This costs time but changes no theorem.

\texttt{NormTargets.PlaneData} includes all coefficients of a monic
equation in $Y$. Companion-matrix powers implement reduction even for
an arbitrary input $Y$-coefficient list. The main general target retains
the essential nonzero-norm hypothesis. The human coprime-degree theorem
proves nonvanishing automatically; the specialized Lean target states
that conclusion but does not yet prove it. These compiled definitions
use rational coefficients and rational points. Generalization to the
coded-field interface of the written theorem is a further formalization
step. The subspace target uses
explicit row-span witnesses and a nonzero annihilator certificate;
Gaussian elimination is the producer described in the human proof.

\subsection{The arithmetic compiler work, in dependency order}

The intended integration uses the pinned repository modules
\path{Syntax.lean}, \path{Deduction.lean}, and \path{PR.lean}
\cite{repo-core}. They already provide scoped terms and formulas,
substitution, intuitionistic deductions, PR function syntax, and
\texttt{PRA.qfInd}. The following sequence avoids a gap between
native correctness and the object-theory claim.
\begin{enumerate}
\item \textbf{Concrete arithmetic.} Introduce PR syntax for signed
integers, positive-denominator rationals, list lookup and folds,
polynomial convolution, bounded maxima and least search. Prove their
defining-equation specifications in \texttt{PRAAx}. Every malformed code
gets a default; correctness predicates assert the valid-input guards.
\item \textbf{Boolean compilation.} Implement the fixed-syntax
translation of Lemma~\ref{lem:booleanization}. Its universal and
existential fold invariants were given explicitly, so this step needs
no induction on an unbounded satisfaction predicate. The compiler can
run in Lean's metatheory; each output must be a genuine derivation.
\item \textbf{Schedule derivations.} Reify the two recursions of
Section~\ref{sec:backward-demands}. Derive the truncated-subtraction
identities and schedule equations from \texttt{PRA.precZero},
\texttt{PRA.precSucc}, equality substitution, and quantifier-free
induction. Build the syntax of the guarded invariant and apply
\texttt{PRA.qfInd}. The native proof gives the exact algebraic template.
\item \textbf{Finite metric algebra.} Prove rational triangle and
finite-minimum specifications by denominator clearing and prefix
induction. Translate \texttt{coverCheck\_correct} and the finite cover
combinators. Formalize the least-neighbor invariant and the explicitly
bounded geometric sum. Then implement the value and witness budgets.
\item \textbf{Rational polynomial analysis.} Prove squared-norm
multiplication, the power-difference identity, coefficient convolution,
synthetic division, and the rational descent lemma. The last is the
substantial inherited obligation. Compose its bounded iterator with
the fixed-length deflation invariant; then substitute $E(t)$ in the
finite product-stability proof and use the adequacy sample indices.
\item \textbf{Finite field algebra.} Prove monic reduction, evaluation
rows, determinant multiplication and row multilinearity, the
Vandermonde calculation, and synthetic removal of linear factors.
For the specialized theorem, formalize the unique-weight argument and
the unique cyclic-shift determinant term. The finite proof has no
irreducibility or divisor-theory dependency to import.
\item \textbf{Output adequacy and transport.} Universally introduce
the numerical parameters only after the finite derivations are built.
For named inputs, retain their validity hypotheses and instantiate them
at the explicitly calculated indices. The existing transport APIs between ordinary iPRA and
the ConciseIPRA scalar calculus can then be applied to an actual arithmetic
derivation; they do not produce that derivation from a native theorem.
\end{enumerate}

For example, in backward transfer the invariant formula is the syntax
of $t\le N\to c(t,h(t;N,K))=1$. The base is a derivation of its
substitution at zero. The step is a derivation in the context containing
that formula and the old context shifted by one variable. Instantiate
the local contract at $h(t+1)$, use the schedule equality, and introduce
the guard implication. After \texttt{PRA.qfInd}, eliminate the universal
quantifier at the term $N$ and rewrite $h(N)=K$. This specifies the
binders, substitutions, and final inference, not merely the shape of
the desired conclusion.

\clearpage
\subsection{Induction inventory}

\begin{longtable}{@{}p{0.23\textwidth}p{0.45\textwidth}p{0.25\textwidth}@{}}
\toprule
Construction & Finite invariant & Terminal resource\\
\midrule
\endhead
Backward transfer & $t\le N\to c(t,h(t))=1$ & $N$ forward steps and
$N$ backward steps.\\[3pt]
Least-neighbor chain & Current point lies in its finite list; accumulated
rational distance is at most the finite geometric sum. & $s+k+3$ as
largest observed list index for membership at $k$.\\[3pt]
Approximate factorization & Monicity and degree; coefficient bounds;
the exact partial telescoping identity; accumulated error. & $n$
divisions, each with its specified descent or grid bound.\\[3pt]
Norm construction & Each matrix column is the specified remainder;
partial determinant sum has its computed coefficient list. & $m^m$
candidate permutation words suffice.\\[3pt]
Norm zero count & Groups partition the input list; removed linear
factors multiply the current quotient to the original norm. & Input
list length and polynomial degrees.\\[3pt]
Subspace fibers & Pivot equations and selected free coordinate;
the annihilator is nonzero and annihilates each supplied row. & Matrix
row and column counts.\\
\bottomrule
\end{longtable}

The native axiom reports are included verbatim in the validation output.
The scheduling results use only \texttt{propext} and \texttt{Quot.sound}.
Pure relational cover composition uses no axioms. The native rational
triangle proof inherits \texttt{Classical.choice} from \texttt{Std}'s
rational-order lemmas. This explains why compiling the native rational
layer is not by itself a certificate of intuitionistic arithmetic
strength. The object development uses coded rational comparisons and
the explicit finite proofs instead.

\section{Repository comparison, validation, and remaining difficulties}
\label{sec:audit}

\subsection{The inspected snapshot}

The source baseline is ProveIt commit
\texttt{052fbaebfc955294b4c02f6532ab7e7a3f760f7a}, retrieved on
15 November 2026 in the user's Pacific-time date convention. The
collection index lists 54 delivered packets and expressly says that
their filing did not rerun, compile, or independently check them
\cite{repo-index}. Our review was targeted. A catalog of packet
descriptions is not an audit of every proof in every packet.

The archive records the selected source paths, pinned URLs, byte sizes,
Git blob identities, and SHA-256 hashes in
\path{provenance/source_manifest.json}. The original FTA article is
included unchanged under \path{dependencies/FTA_iPRA/}; it is the
only substantial inherited existence proof used by the new root
producer. Hash agreement identifies a source; it does not certify its
mathematical correctness.

\clearpage
\subsection{What this adds to the existing reports}

\begin{longtable}{@{}p{0.25\textwidth}p{0.34\textwidth}p{0.34\textwidth}@{}}
\toprule
Earlier result or difficulty & Contribution here & Remaining boundary\\
\midrule
\endhead
Finite contracts and guarded induction; one backwards field-tower
budget in the Platonic-solids packet. & A general nonmonotone backward
transfer rule and a two-recursion compiler for predetermined changing
parameters; scheduling theorems checked in Lean. & A PR demand and
local object derivations must still be supplied.\\[4pt]
Singular Puiseux coefficients cannot be selected uniformly from arbitrary
coefficient names; regular charts already have a derivative-gap method.
& A uniform monic root-support name with explicit coefficient and
Hausdorff budgets, valid through collisions. & No singular Puiseux
coefficient constructor, global branch labels, or multiplicity-matching
theorem.\\[4pt]
Finite observations do not automatically identify one coherent exact
witness. & An explicit aggregate calculus, a least-neighbor point-name
constructor, and approximate optimization over the represented support.
& Point selection is not claimed extensional; approximate minimizers
need not form one exact minimizing name.\\[4pt]
The Siegel route needs a zero-count proof for functions in a finite
evaluation space; the earlier framework handles odd hyperelliptic degree.
& A norm constructor for general monic plane equations, and automatic
sharp degree and zero bounds for every coprime superelliptic pair
$(m,\deg f)$. & The approximation cover, lifting data, and global height
bound are separate mathematical obligations.\\[4pt]
General Mordell--Weil needs actual rational quotient representatives;
Faltings and Roth reports distinguish finite counts from effective
cutoffs. & The framework specifies which finite outputs can be composed
and what the new geometric count supplies. & None of those missing
global producers or cutoffs is obtained here.\\
\bottomrule
\end{longtable}

This comparison makes no mathematical priority claim. Backward budgeting,
completed finite sets, synthetic division, norms, and Vandermonde
determinants are established ideas. The advances relative to this
snapshot are the proved interfaces, explicit PR constructions and
budgets, checked scheduling core, and the broader automatically certified
zero-count family. In particular, the new norm theorem treats coprime
degrees; the preceding power-gap theorem treats a different degree
regime and gives a height conclusion. They are complementary results.

\subsection{What was actually checked}

All five new Lean modules were compiled with Lean 4.32.0. The build
script records the compiler version, return codes, output, source hashes,
and printed axiom dependencies. The distinction between proved native
lemmas and named unproved targets is recorded in
\texttt{proof\_status.json}. No compilation result is presented as a
completed embedded derivation of the main mathematical applications.

The reference programs use integer arithmetic and
\texttt{fractions.Fraction}. The root program checks 12 finite
approximate factorizations, 24 perturbation pairs, six collision cases,
and two rejected invalid hypothesis sets. Its generic finite grid
fallback is specified, but the tests use certified elementary shortcuts
to avoid its enormous worst-case runtime. This validates finite
implementation behavior, not the full descent theorem.

The norm program checks 57 exact degree and leading-coefficient cases,
57 independently constructed companion matrices, 171 independent
determinant evaluations, 12 hyperelliptic specializations, 42
evaluation-row identities, three multiplicity or ramification cases,
the sharp four-point example, and a failed-coprimality example. The
determinants are checked against rational Gaussian elimination,
providing an independent arithmetic computation rather than simply
calling the same determinant routine twice.

The additional aggregate and scheduling reference program checks finite
cover composition, nearest-neighbor chains, nonlinear iteration budgets,
and predetermined parameter-changing recursions. Exact counts and the
reproduction commands are in the validation records. These checks are
useful for finding coding and constant errors; the universal theorems
rest on the written arguments and, where stated, the native Lean proofs.

The PDF was compiled from the delivered sources, its pages were rendered
and visually checked, and the final log was inspected for unresolved
references and layout warnings. Build scripts, code, machine-readable
status, and validation output are included so the result can be
reviewed and extended inside ProveIt.

\subsection{The next mathematical tasks}

The most direct formalization step is the object-level bounded-formula
and guarded-induction compiler. It would serve all three applications
and many existing packets. The deepest analytic dependency in this
package is the rational descent proof; formalizing it would turn the
root-support construction into a reusable deeply embedded theorem.
The determinant and weighted-degree argument is a separate finite
algebra project with no analytic dependency.

For the larger Diophantine program, the new subspace-fiber count now
supplies the geometric bound for the stated superelliptic evaluation
spaces. The finite-tuple approximation cover remains the next genuinely
new arithmetic ingredient. A cardinality bound obtained after that step
would still have a different logical form from a positive coordinate
height bound. For general rational-group generators, an actual
quotient-covering list remains the missing input identified in the
Mordell--Weil packet \cite{repo-mordell}. The unresolved positive
height targets in the Faltings packet remain unresolved
\cite{repo-faltings}; this report draws no independence or unprovability
conclusion from that fact.

\clearpage
\begin{thebibliography}{99}

\small

\bibitem{richman2000} Fred Richman, \emph{The fundamental theorem of algebra: a constructive development without choice}, Pacific Journal of Mathematics \textbf{196} (2000), no.~1, 213--230. DOI: 10.2140/pjm.2000.196.213. \url{https://msp.org/pjm/2000/196-1/pjm-v196-n1-p10-p.pdf}.

\medskip\noindent The following repository references are pinned to ProveIt commit \texttt{052fbaebfc955294b4c02f6532ab7e7a3f760f7a}. Titles identify the selected packets; their statements are not assumed independently verified merely because they are present in the repository.

\bibitem{repo-index} ProveIt contributors, \emph{Delivered iPRA proof packets}.\par\nopagebreak\href{https://github.com/VladimirReshetnikov/ProveIt/blob/052fbaebfc955294b4c02f6532ab7e7a3f760f7a/Logic/PrimitiveRecursiveArithmetic/docs/proofs/README.md}{\nolinkurl{Logic/PrimitiveRecursiveArithmetic/docs/proofs/README.md}}.

\bibitem{repo-core} ProveIt contributors, \emph{Syntax and deduction of intuitionistic primitive recursive arithmetic}.\par\nopagebreak\href{https://github.com/VladimirReshetnikov/ProveIt/blob/052fbaebfc955294b4c02f6532ab7e7a3f760f7a/Logic/PrimitiveRecursiveArithmetic/ConciseIPRA/Lean/ConciseIPRA/PR.lean}{\nolinkurl{Logic/PrimitiveRecursiveArithmetic/ConciseIPRA/Lean/ConciseIPRA/PR.lean}}.

\bibitem{repo-contracts} ProveIt contributors, \emph{Finite contracts for mathematics in intuitionistic primitive recursive arithmetic}.\par\nopagebreak\href{https://github.com/VladimirReshetnikov/ProveIt/blob/052fbaebfc955294b4c02f6532ab7e7a3f760f7a/Logic/PrimitiveRecursiveArithmetic/docs/proofs/foundations/ipra_framework/paper/ipra_framework.tex}{\nolinkurl{Logic/PrimitiveRecursiveArithmetic/docs/proofs/foundations/ipra_framework/paper/ipra_framework.tex}}.

\bibitem{repo-embedding} ProveIt contributors, \emph{Mathematics in iPRA through finite observations}.\par\nopagebreak\href{https://github.com/VladimirReshetnikov/ProveIt/blob/052fbaebfc955294b4c02f6532ab7e7a3f760f7a/Logic/PrimitiveRecursiveArithmetic/docs/proofs/foundations/ipra_embedding_framework/report.tex}{\nolinkurl{Logic/PrimitiveRecursiveArithmetic/docs/proofs/foundations/ipra_embedding_framework/report.tex}}.

\bibitem{repo-first-framework} ProveIt contributors, \emph{iPRA certificate-and-modulus framework}.\par\nopagebreak\href{https://github.com/VladimirReshetnikov/ProveIt/blob/052fbaebfc955294b4c02f6532ab7e7a3f760f7a/Logic/PrimitiveRecursiveArithmetic/docs/proofs/foundations/ipra_framework_2026-11-15/framework.tex}{\nolinkurl{Logic/PrimitiveRecursiveArithmetic/docs/proofs/foundations/ipra_framework_2026-11-15/framework.tex}}.

\bibitem{repo-platonic} ProveIt contributors, \emph{Platonic solids: arithmetic formalization and backward coefficient budgets}.\par\nopagebreak\href{https://github.com/VladimirReshetnikov/ProveIt/blob/052fbaebfc955294b4c02f6532ab7e7a3f760f7a/Logic/PrimitiveRecursiveArithmetic/docs/proofs/geometry/platonic_solids_ipra/paper/sections/05_formalization.tex}{\nolinkurl{Logic/PrimitiveRecursiveArithmetic/docs/proofs/geometry/platonic_solids_ipra/paper/sections/05_formalization.tex}}.

\bibitem{repo-fta} ProveIt contributors, \emph{Fundamental theorem of algebra in intuitionistic PRA}.\par\nopagebreak\href{https://github.com/VladimirReshetnikov/ProveIt/blob/052fbaebfc955294b4c02f6532ab7e7a3f760f7a/Logic/PrimitiveRecursiveArithmetic/docs/proofs/algebra/FTA_iPRA/FTA_iPRA.tex}{\nolinkurl{Logic/PrimitiveRecursiveArithmetic/docs/proofs/algebra/FTA_iPRA/FTA_iPRA.tex}}.

\bibitem{repo-puiseux} ProveIt contributors, \emph{Puiseux expansions in intuitionistic primitive recursive arithmetic}.\par\nopagebreak\href{https://github.com/VladimirReshetnikov/ProveIt/blob/052fbaebfc955294b4c02f6532ab7e7a3f760f7a/Logic/PrimitiveRecursiveArithmetic/docs/proofs/algebra/puiseux_ipra/puiseux_ipra.tex}{\nolinkurl{Logic/PrimitiveRecursiveArithmetic/docs/proofs/algebra/puiseux_ipra/puiseux_ipra.tex}}.

\bibitem{repo-cayley} ProveIt contributors, \emph{Cayley--Hamilton theorem in iPRA}.\par\nopagebreak\href{https://github.com/VladimirReshetnikov/ProveIt/blob/052fbaebfc955294b4c02f6532ab7e7a3f760f7a/Logic/PrimitiveRecursiveArithmetic/docs/proofs/algebra/cayley_hamilton_ipra/README.md}{\nolinkurl{Logic/PrimitiveRecursiveArithmetic/docs/proofs/algebra/cayley_hamilton_ipra/README.md}}.

\bibitem{repo-siegel} ProveIt contributors, \emph{Siegel theorem packet: the general route and remaining obligations}.\par\nopagebreak\href{https://github.com/VladimirReshetnikov/ProveIt/blob/052fbaebfc955294b4c02f6532ab7e7a3f760f7a/Logic/PrimitiveRecursiveArithmetic/docs/proofs/diophantine-number-theory/siegel_ipra/paper/05_general_route.tex}{\nolinkurl{Logic/PrimitiveRecursiveArithmetic/docs/proofs/diophantine-number-theory/siegel_ipra/paper/05_general_route.tex}}.

\bibitem{repo-mordell} ProveIt contributors, \emph{Constructive Mordell--Weil descent in iPRA}.\par\nopagebreak\href{https://github.com/VladimirReshetnikov/ProveIt/blob/052fbaebfc955294b4c02f6532ab7e7a3f760f7a/Logic/PrimitiveRecursiveArithmetic/docs/proofs/diophantine-number-theory/mordell_weil_ipra/README.txt}{\nolinkurl{Logic/PrimitiveRecursiveArithmetic/docs/proofs/diophantine-number-theory/mordell_weil_ipra/README.txt}}.

\bibitem{repo-faltings} ProveIt contributors, \emph{Faltings theorem and intuitionistic PRA}.\par\nopagebreak\href{https://github.com/VladimirReshetnikov/ProveIt/blob/052fbaebfc955294b4c02f6532ab7e7a3f760f7a/Logic/PrimitiveRecursiveArithmetic/docs/proofs/diophantine-number-theory/faltings_ipra/README.md}{\nolinkurl{Logic/PrimitiveRecursiveArithmetic/docs/proofs/diophantine-number-theory/faltings_ipra/README.md}}.

\end{thebibliography}

\end{document}
